% Hybrid Linear-Convex Confidence (HLCC) Loss for Calibrated Language Models
% Compile with: pdflatex hlcc_paper.tex && bibtex hlcc_paper && pdflatex hlcc_paper && pdflatex hlcc_paper

\documentclass[11pt,a4paper]{article}

% Packages
\usepackage[utf8]{inputenc}
\usepackage[T1]{fontenc}
\usepackage{amsmath,amssymb,amsthm}
\usepackage{booktabs}
\usepackage{graphicx}
\usepackage[numbers]{natbib}
\usepackage{hyperref}
\usepackage{cleveref}
\usepackage{algorithm}
\usepackage{algorithmic}
\usepackage{multirow}
\usepackage{subcaption}
\usepackage[margin=2.5cm]{geometry}
\usepackage{enumitem}
\usepackage{xcolor}
\usepackage{float}
\usepackage{csvsimple-l3}

% Theorem environments
\newtheorem{theorem}{Theorem}
\newtheorem{proposition}{Proposition}
\newtheorem{corollary}{Corollary}
\newtheorem{definition}{Definition}
\newtheorem{remark}{Remark}

% Custom commands
\newcommand{\hlcc}{\textsc{hlcc}}
\newcommand{\E}{\mathbb{E}}
\newcommand{\R}{\mathbb{R}}
\newcommand{\cstar}{c^{*}}
\DeclareMathOperator*{\argmax}{arg\,max}
\DeclareMathOperator*{\argmin}{arg\,min}

\title{Hybrid Linear-Convex Confidence Loss\\for Calibrated Language Models}
\author{
  Simon [Last Name]\\
  \texttt{[email]}
}
\date{\today}

\begin{document}
\maketitle

% Load all data-driven numbers (generated by generate_paper_tables.py)
% AUTO-GENERATED by generate_paper_tables.py
% DO NOT EDIT MANUALLY -- re-run the script to update.
% Generated: 2026-02-25T07:29:23.784395

% ============================================
% HLCC Scoring Function Constants
% SOURCE: D:/git/NNConfidence\paper_constants.yaml
% ============================================
\newcommand{\hlccCorrectCoeff}{1}
\newcommand{\hlccIncorrectCoeff}{-2}
\newcommand{\hlccCrossover}{0.75}
\newcommand{\hlccSaturation}{0.8}
\newcommand{\hlccMarginalRatio}{4}
\newcommand{\hlccSafeGuess}{0.25}

% ============================================
% Training Hyperparameters
% SOURCE: D:/git/NNConfidence\paper_constants.yaml
% ============================================
\newcommand{\expEpochs}{25}
\newcommand{\expTrainExamples}{900}
\newcommand{\expValExamples}{100}
\newcommand{\expBatchSize}{4}
\newcommand{\expLR}{1 \times 10^{-4}}
\newcommand{\expWeightDecay}{0.01}
\newcommand{\expAlpha}{0.5}
\newcommand{\expWarmupSteps}{1000}
\newcommand{\expExamplesPerDataset}{300}

% ============================================
% Calibration Parameters
% SOURCE: D:/git/NNConfidence\paper_constants.yaml
% ============================================
\newcommand{\expCalibrationExamples}{300}
\newcommand{\expSelfConSamples}{10}
\newcommand{\expPolygraphExamples}{817}
\newcommand{\expPolygraphRetrainExamples}{600}
\newcommand{\expPolygraphRetrainEpochs}{15}
\newcommand{\expLayerExamples}{90}

% ============================================
% Hardware
% SOURCE: D:/git/NNConfidence\paper_constants.yaml
% ============================================
\newcommand{\hwGPU}{RTX PRO 4500 Blackwell}
\newcommand{\hwVRAM}{32}
\newcommand{\hwCPU}{AMD Ryzen 9}
\newcommand{\hwRAM}{64}

% ============================================
% NormShift Ablation Results (from checkpoints)
% SOURCE: D:/git/NNConfidence\data/checkpoints/<model>_norm_shift/best_norm_shift_*.pt
% ============================================
% Mistral-7B combined (epoch 21)
% SOURCE: D:/git/NNConfidence\data/checkpoints/mistral-7b_norm_shift/best_norm_shift_combined.pt
\newcommand{\mistralCombinedEpoch}{21}
\newcommand{\mistralCombinedAcc}{0.53}
\newcommand{\mistralCombinedECE}{0.020}
\newcommand{\mistralCombinedBrier}{0.020}
\newcommand{\mistralCombinedHLCC}{1.039}
\newcommand{\mistralCombinedConf}{0.510}
\newcommand{\mistralCombinedCalGap}{0.020}
% Mistral-7B norm_shift_only (epoch 0)
% SOURCE: D:/git/NNConfidence\data/checkpoints/mistral-7b_norm_shift/best_norm_shift_norm_shift_only.pt
\newcommand{\mistralNSOEpoch}{0}
\newcommand{\mistralNSOAcc}{0.53}
\newcommand{\mistralNSOECE}{0.054}
\newcommand{\mistralNSOBrier}{0.252}
\newcommand{\mistralNSOHLCC}{0.519}
\newcommand{\mistralNSOConf}{0.584}
\newcommand{\mistralNSOCalGap}{0.054}

% Qwen2.5-7B combined (epoch 19)
% SOURCE: D:/git/NNConfidence\data/checkpoints/qwen2.5-7b_norm_shift/best_norm_shift_combined.pt
\newcommand{\qwenCombinedEpoch}{19}
\newcommand{\qwenCombinedAcc}{0.32}
\newcommand{\qwenCombinedECE}{0.060}
\newcommand{\qwenCombinedBrier}{0.060}
\newcommand{\qwenCombinedHLCC}{0.540}
\newcommand{\qwenCombinedConf}{0.340}
\newcommand{\qwenCombinedCalGap}{0.020}
% Qwen2.5-7B norm_shift_only (epoch 20)
% SOURCE: D:/git/NNConfidence\data/checkpoints/qwen2.5-7b_norm_shift/best_norm_shift_norm_shift_only.pt
\newcommand{\qwenNSOEpoch}{20}
\newcommand{\qwenNSOAcc}{0.32}
\newcommand{\qwenNSOECE}{0.133}
\newcommand{\qwenNSOBrier}{0.236}
\newcommand{\qwenNSOHLCC}{0.332}
\newcommand{\qwenNSOConf}{0.187}
\newcommand{\qwenNSOCalGap}{0.133}

% Llama-3.1-8B combined (epoch 22)
% SOURCE: D:/git/NNConfidence\data/checkpoints/llama3.1-8b_norm_shift/best_norm_shift_combined.pt
\newcommand{\llamaCombinedEpoch}{22}
\newcommand{\llamaCombinedAcc}{0.92}
\newcommand{\llamaCombinedECE}{0.040}
\newcommand{\llamaCombinedBrier}{0.040}
\newcommand{\llamaCombinedHLCC}{1.760}
\newcommand{\llamaCombinedConf}{0.960}
\newcommand{\llamaCombinedCalGap}{0.040}
% Llama-3.1-8B norm_shift_only (epoch 0)
% SOURCE: D:/git/NNConfidence\data/checkpoints/llama3.1-8b_norm_shift/best_norm_shift_norm_shift_only.pt
\newcommand{\llamaNSOEpoch}{0}
\newcommand{\llamaNSOAcc}{0.92}
\newcommand{\llamaNSOECE}{0.517}
\newcommand{\llamaNSOBrier}{0.341}
\newcommand{\llamaNSOHLCC}{1.265}
\newcommand{\llamaNSOConf}{0.403}
\newcommand{\llamaNSOCalGap}{0.517}

% Qwen2.5-14B combined (epoch 23)
% SOURCE: D:/git/NNConfidence\data/checkpoints/qwen2.5-14b_norm_shift/best_norm_shift_combined.pt
\newcommand{\qwenLCombinedEpoch}{23}
\newcommand{\qwenLCombinedAcc}{0.89}
\newcommand{\qwenLCombinedECE}{0.072}
\newcommand{\qwenLCombinedBrier}{0.063}
\newcommand{\qwenLCombinedHLCC}{1.657}
\newcommand{\qwenLCombinedConf}{0.927}
\newcommand{\qwenLCombinedCalGap}{0.037}
% Qwen2.5-14B norm_shift_only (epoch 20)
% SOURCE: D:/git/NNConfidence\data/checkpoints/qwen2.5-14b_norm_shift/best_norm_shift_norm_shift_only.pt
\newcommand{\qwenLNSOEpoch}{20}
\newcommand{\qwenLNSOAcc}{0.89}
\newcommand{\qwenLNSOECE}{0.164}
\newcommand{\qwenLNSOBrier}{0.138}
\newcommand{\qwenLNSOHLCC}{1.429}
\newcommand{\qwenLNSOConf}{0.726}
\newcommand{\qwenLNSOCalGap}{0.164}

% DeepSeek-R1-14B combined (epoch 22)
% SOURCE: D:/git/NNConfidence\data/checkpoints/deepseek-r1-14b_norm_shift/best_norm_shift_combined.pt
\newcommand{\deepseekCombinedEpoch}{22}
\newcommand{\deepseekCombinedAcc}{0.48}
\newcommand{\deepseekCombinedECE}{0.108}
\newcommand{\deepseekCombinedBrier}{0.105}
\newcommand{\deepseekCombinedHLCC}{0.781}
\newcommand{\deepseekCombinedConf}{0.501}
\newcommand{\deepseekCombinedCalGap}{0.021}
% DeepSeek-R1-14B norm_shift_only (epoch 0)
% SOURCE: D:/git/NNConfidence\data/checkpoints/deepseek-r1-14b_norm_shift/best_norm_shift_norm_shift_only.pt
\newcommand{\deepseekNSOEpoch}{0}
\newcommand{\deepseekNSOAcc}{0.48}
\newcommand{\deepseekNSOECE}{0.442}
\newcommand{\deepseekNSOBrier}{0.445}
\newcommand{\deepseekNSOHLCC}{0.038}
\newcommand{\deepseekNSOConf}{0.922}
\newcommand{\deepseekNSOCalGap}{0.442}


% ============================================
% Calibration Baseline Results
% SOURCE: D:/git/NNConfidence\data/results/calibration_progress_*_truthfulqa.json
% ============================================
% Mistral-7B (14 methods, 300 examples)
\newcommand{\mistralCalNExamples}{300}
\newcommand{\mistralCalAcc}{0.467}
\newcommand{\mistralEntropyAcc}{0.467}
\newcommand{\mistralEntropyConf}{0.855}
\newcommand{\mistralEntropyCalGap}{0.389}
\newcommand{\mistralEntropyECE}{0.389}
\newcommand{\mistralEntropyBrier}{0.393}
\newcommand{\mistralEntropyHLCC}{0.105}
\newcommand{\mistralEntropyAUROC}{0.602}
\newcommand{\mistralEntropyAUPRC}{0.531}
\newcommand{\mistralEntropyPRR}{0.159}
\newcommand{\mistralEntropyAccSeventy}{0.505}
\newcommand{\mistralEntropyAccEighty}{0.492}
\newcommand{\mistralEntropyAccNinety}{0.481}
\newcommand{\mistralTopkAcc}{0.467}
\newcommand{\mistralTopkConf}{0.891}
\newcommand{\mistralTopkCalGap}{0.424}
\newcommand{\mistralTopkECE}{0.428}
\newcommand{\mistralTopkBrier}{0.422}
\newcommand{\mistralTopkHLCC}{0.060}
\newcommand{\mistralTopkAUROC}{0.602}
\newcommand{\mistralTopkAUPRC}{0.545}
\newcommand{\mistralTopkPRR}{0.194}
\newcommand{\mistralTopkAccSeventy}{0.505}
\newcommand{\mistralTopkAccEighty}{0.492}
\newcommand{\mistralTopkAccNinety}{0.485}
\newcommand{\mistralMSPAcc}{0.467}
\newcommand{\mistralMSPConf}{0.819}
\newcommand{\mistralMSPCalGap}{0.353}
\newcommand{\mistralMSPECE}{0.353}
\newcommand{\mistralMSPBrier}{0.352}
\newcommand{\mistralMSPHLCC}{0.214}
\newcommand{\mistralMSPAUROC}{0.711}
\newcommand{\mistralMSPAUPRC}{0.706}
\newcommand{\mistralMSPPRR}{0.513}
\newcommand{\mistralMSPAccSeventy}{0.548}
\newcommand{\mistralMSPAccEighty}{0.512}
\newcommand{\mistralMSPAccNinety}{0.496}
\newcommand{\mistralGNLLAcc}{0.467}
\newcommand{\mistralGNLLConf}{0.819}
\newcommand{\mistralGNLLCalGap}{0.353}
\newcommand{\mistralGNLLECE}{0.353}
\newcommand{\mistralGNLLBrier}{0.352}
\newcommand{\mistralGNLLHLCC}{0.214}
\newcommand{\mistralGNLLAUROC}{0.711}
\newcommand{\mistralGNLLAUPRC}{0.706}
\newcommand{\mistralGNLLPRR}{0.513}
\newcommand{\mistralGNLLAccSeventy}{0.548}
\newcommand{\mistralGNLLAccEighty}{0.512}
\newcommand{\mistralGNLLAccNinety}{0.496}
\newcommand{\mistralVerbAcc}{0.467}
\newcommand{\mistralVerbConf}{0.936}
\newcommand{\mistralVerbCalGap}{0.469}
\newcommand{\mistralVerbECE}{0.502}
\newcommand{\mistralVerbBrier}{0.516}
\newcommand{\mistralVerbHLCC}{-0.075}
\newcommand{\mistralVerbAUROC}{0.595}
\newcommand{\mistralVerbAUPRC}{0.557}
\newcommand{\mistralVerbPRR}{0.222}
\newcommand{\mistralVerbAccSeventy}{0.543}
\newcommand{\mistralVerbAccEighty}{0.504}
\newcommand{\mistralVerbAccNinety}{0.463}
\newcommand{\mistralSelfConAcc}{0.290}
\newcommand{\mistralSelfConConf}{0.712}
\newcommand{\mistralSelfConCalGap}{0.422}
\newcommand{\mistralSelfConECE}{0.422}
\newcommand{\mistralSelfConBrier}{0.410}
\newcommand{\mistralSelfConHLCC}{-0.254}
\newcommand{\mistralSelfConAUROC}{0.556}
\newcommand{\mistralSelfConAUPRC}{0.338}
\newcommand{\mistralSelfConPRR}{0.114}
\newcommand{\mistralSelfConAccSeventy}{0.314}
\newcommand{\mistralSelfConAccEighty}{0.300}
\newcommand{\mistralSelfConAccNinety}{0.296}
\newcommand{\mistralTempScaleAcc}{0.467}
\newcommand{\mistralTempScaleConf}{0.982}
\newcommand{\mistralTempScaleCalGap}{0.515}
\newcommand{\mistralTempScaleECE}{0.515}
\newcommand{\mistralTempScaleBrier}{0.510}
\newcommand{\mistralTempScaleHLCC}{-0.090}
\newcommand{\mistralTempScaleAUROC}{0.651}
\newcommand{\mistralTempScaleAUPRC}{0.641}
\newcommand{\mistralTempScalePRR}{0.393}
\newcommand{\mistralTempScaleAccSeventy}{0.533}
\newcommand{\mistralTempScaleAccEighty}{0.521}
\newcommand{\mistralTempScaleAccNinety}{0.489}
\newcommand{\mistralIsotonicAcc}{0.467}
\newcommand{\mistralIsotonicConf}{0.467}
\newcommand{\mistralIsotonicCalGap}{0.000}
\newcommand{\mistralIsotonicECE}{0.000}
\newcommand{\mistralIsotonicBrier}{0.230}
\newcommand{\mistralIsotonicHLCC}{0.485}
\newcommand{\mistralIsotonicAUROC}{0.654}
\newcommand{\mistralIsotonicAUPRC}{0.627}
\newcommand{\mistralIsotonicPRR}{0.368}
\newcommand{\mistralIsotonicAccSeventy}{0.529}
\newcommand{\mistralIsotonicAccEighty}{0.496}
\newcommand{\mistralIsotonicAccNinety}{0.481}
\newcommand{\mistralPlattAcc}{0.467}
\newcommand{\mistralPlattConf}{0.467}
\newcommand{\mistralPlattCalGap}{0.000}
\newcommand{\mistralPlattECE}{0.050}
\newcommand{\mistralPlattBrier}{0.243}
\newcommand{\mistralPlattHLCC}{0.462}
\newcommand{\mistralPlattAUROC}{0.602}
\newcommand{\mistralPlattAUPRC}{0.531}
\newcommand{\mistralPlattPRR}{0.159}
\newcommand{\mistralPlattAccSeventy}{0.505}
\newcommand{\mistralPlattAccEighty}{0.492}
\newcommand{\mistralPlattAccNinety}{0.481}
\newcommand{\mistralNormShiftAcc}{0.467}
\newcommand{\mistralNormShiftConf}{0.497}
\newcommand{\mistralNormShiftCalGap}{0.030}
\newcommand{\mistralNormShiftECE}{0.030}
\newcommand{\mistralNormShiftBrier}{0.025}
\newcommand{\mistralNormShiftHLCC}{0.884}
\newcommand{\mistralNormShiftAUROC}{0.999}
\newcommand{\mistralNormShiftAUPRC}{0.999}
\newcommand{\mistralNormShiftPRR}{0.999}
\newcommand{\mistralNormShiftAccSeventy}{0.667}
\newcommand{\mistralNormShiftAccEighty}{0.583}
\newcommand{\mistralNormShiftAccNinety}{0.519}
\newcommand{\mistralSemEntAcc}{0.300}
\newcommand{\mistralSemEntConf}{0.572}
\newcommand{\mistralSemEntCalGap}{0.272}
\newcommand{\mistralSemEntECE}{0.328}
\newcommand{\mistralSemEntBrier}{0.354}
\newcommand{\mistralSemEntHLCC}{-0.076}
\newcommand{\mistralSemEntAUROC}{0.518}
\newcommand{\mistralSemEntAUPRC}{0.323}
\newcommand{\mistralSemEntPRR}{0.054}
\newcommand{\mistralSemEntAccSeventy}{0.295}
\newcommand{\mistralSemEntAccEighty}{0.296}
\newcommand{\mistralSemEntAccNinety}{0.289}
\newcommand{\mistralCotConfAcc}{0.513}
\newcommand{\mistralCotConfConf}{0.956}
\newcommand{\mistralCotConfCalGap}{0.443}
\newcommand{\mistralCotConfECE}{0.463}
\newcommand{\mistralCotConfBrier}{0.449}
\newcommand{\mistralCotConfHLCC}{0.124}
\newcommand{\mistralCotConfAUROC}{0.581}
\newcommand{\mistralCotConfAUPRC}{0.601}
\newcommand{\mistralCotConfPRR}{0.214}
\newcommand{\mistralCotConfAccSeventy}{0.571}
\newcommand{\mistralCotConfAccEighty}{0.550}
\newcommand{\mistralCotConfAccNinety}{0.533}
\newcommand{\mistralLVUAcc}{0.467}
\newcommand{\mistralLVUConf}{0.463}
\newcommand{\mistralLVUCalGap}{0.004}
\newcommand{\mistralLVUECE}{0.122}
\newcommand{\mistralLVUBrier}{0.287}
\newcommand{\mistralLVUHLCC}{0.412}
\newcommand{\mistralLVUAUROC}{0.464}
\newcommand{\mistralLVUAUPRC}{0.453}
\newcommand{\mistralLVUPRR}{-0.044}
\newcommand{\mistralLVUAccSeventy}{0.443}
\newcommand{\mistralLVUAccEighty}{0.446}
\newcommand{\mistralLVUAccNinety}{0.441}
\newcommand{\mistralCoCoAAcc}{0.277}
\newcommand{\mistralCoCoAConf}{0.783}
\newcommand{\mistralCoCoACalGap}{0.506}
\newcommand{\mistralCoCoAECE}{0.506}
\newcommand{\mistralCoCoABrier}{0.456}
\newcommand{\mistralCoCoAHLCC}{-0.385}
\newcommand{\mistralCoCoAAUROC}{0.585}
\newcommand{\mistralCoCoAAUPRC}{0.323}
\newcommand{\mistralCoCoAPRR}{0.115}
\newcommand{\mistralCoCoAAccSeventy}{0.319}
\newcommand{\mistralCoCoAAccEighty}{0.312}
\newcommand{\mistralCoCoAAccNinety}{0.300}

% Qwen2.5-7B (14 methods, 300 examples)
\newcommand{\qwenCalNExamples}{300}
\newcommand{\qwenCalAcc}{0.493}
\newcommand{\qwenEntropyAcc}{0.493}
\newcommand{\qwenEntropyConf}{0.942}
\newcommand{\qwenEntropyCalGap}{0.449}
\newcommand{\qwenEntropyECE}{0.449}
\newcommand{\qwenEntropyBrier}{0.425}
\newcommand{\qwenEntropyHLCC}{0.125}
\newcommand{\qwenEntropyAUROC}{0.899}
\newcommand{\qwenEntropyAUPRC}{0.895}
\newcommand{\qwenEntropyPRR}{0.834}
\newcommand{\qwenEntropyAccSeventy}{0.671}
\newcommand{\qwenEntropyAccEighty}{0.608}
\newcommand{\qwenEntropyAccNinety}{0.548}
\newcommand{\qwenTopkAcc}{0.493}
\newcommand{\qwenTopkConf}{0.970}
\newcommand{\qwenTopkCalGap}{0.476}
\newcommand{\qwenTopkECE}{0.476}
\newcommand{\qwenTopkBrier}{0.465}
\newcommand{\qwenTopkHLCC}{0.050}
\newcommand{\qwenTopkAUROC}{0.757}
\newcommand{\qwenTopkAUPRC}{0.769}
\newcommand{\qwenTopkPRR}{0.601}
\newcommand{\qwenTopkAccSeventy}{0.605}
\newcommand{\qwenTopkAccEighty}{0.567}
\newcommand{\qwenTopkAccNinety}{0.537}
\newcommand{\qwenMSPAcc}{0.493}
\newcommand{\qwenMSPConf}{0.770}
\newcommand{\qwenMSPCalGap}{0.277}
\newcommand{\qwenMSPECE}{0.277}
\newcommand{\qwenMSPBrier}{0.310}
\newcommand{\qwenMSPHLCC}{0.340}
\newcommand{\qwenMSPAUROC}{0.657}
\newcommand{\qwenMSPAUPRC}{0.589}
\newcommand{\qwenMSPPRR}{0.228}
\newcommand{\qwenMSPAccSeventy}{0.552}
\newcommand{\qwenMSPAccEighty}{0.546}
\newcommand{\qwenMSPAccNinety}{0.519}
\newcommand{\qwenGNLLAcc}{0.493}
\newcommand{\qwenGNLLConf}{0.770}
\newcommand{\qwenGNLLCalGap}{0.277}
\newcommand{\qwenGNLLECE}{0.277}
\newcommand{\qwenGNLLBrier}{0.310}
\newcommand{\qwenGNLLHLCC}{0.340}
\newcommand{\qwenGNLLAUROC}{0.657}
\newcommand{\qwenGNLLAUPRC}{0.589}
\newcommand{\qwenGNLLPRR}{0.228}
\newcommand{\qwenGNLLAccSeventy}{0.552}
\newcommand{\qwenGNLLAccEighty}{0.546}
\newcommand{\qwenGNLLAccNinety}{0.519}
\newcommand{\qwenVerbAcc}{0.493}
\newcommand{\qwenVerbConf}{0.328}
\newcommand{\qwenVerbCalGap}{0.165}
\newcommand{\qwenVerbECE}{0.311}
\newcommand{\qwenVerbBrier}{0.404}
\newcommand{\qwenVerbHLCC}{0.399}
\newcommand{\qwenVerbAUROC}{0.384}
\newcommand{\qwenVerbAUPRC}{0.419}
\newcommand{\qwenVerbPRR}{-0.253}
\newcommand{\qwenVerbAccSeventy}{0.462}
\newcommand{\qwenVerbAccEighty}{0.479}
\newcommand{\qwenVerbAccNinety}{0.485}
\newcommand{\qwenSelfConAcc}{0.683}
\newcommand{\qwenSelfConConf}{0.949}
\newcommand{\qwenSelfConCalGap}{0.265}
\newcommand{\qwenSelfConECE}{0.269}
\newcommand{\qwenSelfConBrier}{0.268}
\newcommand{\qwenSelfConHLCC}{0.825}
\newcommand{\qwenSelfConAUROC}{0.626}
\newcommand{\qwenSelfConAUPRC}{0.754}
\newcommand{\qwenSelfConPRR}{0.254}
\newcommand{\qwenSelfConAccSeventy}{0.738}
\newcommand{\qwenSelfConAccEighty}{0.742}
\newcommand{\qwenSelfConAccNinety}{0.726}
\newcommand{\qwenTempScaleAcc}{0.493}
\newcommand{\qwenTempScaleConf}{0.980}
\newcommand{\qwenTempScaleCalGap}{0.487}
\newcommand{\qwenTempScaleECE}{0.487}
\newcommand{\qwenTempScaleBrier}{0.487}
\newcommand{\qwenTempScaleHLCC}{0.009}
\newcommand{\qwenTempScaleAUROC}{0.652}
\newcommand{\qwenTempScaleAUPRC}{0.586}
\newcommand{\qwenTempScalePRR}{0.235}
\newcommand{\qwenTempScaleAccSeventy}{0.571}
\newcommand{\qwenTempScaleAccEighty}{0.538}
\newcommand{\qwenTempScaleAccNinety}{0.500}
\newcommand{\qwenIsotonicAcc}{0.493}
\newcommand{\qwenIsotonicConf}{0.493}
\newcommand{\qwenIsotonicCalGap}{0.000}
\newcommand{\qwenIsotonicECE}{0.000}
\newcommand{\qwenIsotonicBrier}{0.119}
\newcommand{\qwenIsotonicHLCC}{0.746}
\newcommand{\qwenIsotonicAUROC}{0.909}
\newcommand{\qwenIsotonicAUPRC}{0.912}
\newcommand{\qwenIsotonicPRR}{0.860}
\newcommand{\qwenIsotonicAccSeventy}{0.676}
\newcommand{\qwenIsotonicAccEighty}{0.608}
\newcommand{\qwenIsotonicAccNinety}{0.548}
\newcommand{\qwenPlattAcc}{0.493}
\newcommand{\qwenPlattConf}{0.493}
\newcommand{\qwenPlattCalGap}{0.000}
\newcommand{\qwenPlattECE}{0.037}
\newcommand{\qwenPlattBrier}{0.130}
\newcommand{\qwenPlattHLCC}{0.731}
\newcommand{\qwenPlattAUROC}{0.899}
\newcommand{\qwenPlattAUPRC}{0.895}
\newcommand{\qwenPlattPRR}{0.834}
\newcommand{\qwenPlattAccSeventy}{0.671}
\newcommand{\qwenPlattAccEighty}{0.608}
\newcommand{\qwenPlattAccNinety}{0.548}
\newcommand{\qwenNormShiftAcc}{0.493}
\newcommand{\qwenNormShiftConf}{0.450}
\newcommand{\qwenNormShiftCalGap}{0.043}
\newcommand{\qwenNormShiftECE}{0.305}
\newcommand{\qwenNormShiftBrier}{0.339}
\newcommand{\qwenNormShiftHLCC}{0.460}
\newcommand{\qwenNormShiftAUROC}{0.633}
\newcommand{\qwenNormShiftAUPRC}{0.574}
\newcommand{\qwenNormShiftPRR}{0.211}
\newcommand{\qwenNormShiftAccSeventy}{0.571}
\newcommand{\qwenNormShiftAccEighty}{0.558}
\newcommand{\qwenNormShiftAccNinety}{0.519}
\newcommand{\qwenSemEntAcc}{0.687}
\newcommand{\qwenSemEntConf}{0.933}
\newcommand{\qwenSemEntCalGap}{0.247}
\newcommand{\qwenSemEntECE}{0.261}
\newcommand{\qwenSemEntBrier}{0.277}
\newcommand{\qwenSemEntHLCC}{0.824}
\newcommand{\qwenSemEntAUROC}{0.602}
\newcommand{\qwenSemEntAUPRC}{0.739}
\newcommand{\qwenSemEntPRR}{0.193}
\newcommand{\qwenSemEntAccSeventy}{0.729}
\newcommand{\qwenSemEntAccEighty}{0.733}
\newcommand{\qwenSemEntAccNinety}{0.711}
\newcommand{\qwenCotConfAcc}{0.657}
\newcommand{\qwenCotConfConf}{0.878}
\newcommand{\qwenCotConfCalGap}{0.222}
\newcommand{\qwenCotConfECE}{0.262}
\newcommand{\qwenCotConfBrier}{0.277}
\newcommand{\qwenCotConfHLCC}{0.726}
\newcommand{\qwenCotConfAUROC}{0.637}
\newcommand{\qwenCotConfAUPRC}{0.752}
\newcommand{\qwenCotConfPRR}{0.312}
\newcommand{\qwenCotConfAccSeventy}{0.700}
\newcommand{\qwenCotConfAccEighty}{0.679}
\newcommand{\qwenCotConfAccNinety}{0.663}
\newcommand{\qwenLVUAcc}{0.493}
\newcommand{\qwenLVUConf}{0.369}
\newcommand{\qwenLVUCalGap}{0.124}
\newcommand{\qwenLVUECE}{0.130}
\newcommand{\qwenLVUBrier}{0.285}
\newcommand{\qwenLVUHLCC}{0.518}
\newcommand{\qwenLVUAUROC}{0.554}
\newcommand{\qwenLVUAUPRC}{0.510}
\newcommand{\qwenLVUPRR}{0.037}
\newcommand{\qwenLVUAccSeventy}{0.529}
\newcommand{\qwenLVUAccEighty}{0.521}
\newcommand{\qwenLVUAccNinety}{0.507}
\newcommand{\qwenCoCoAAcc}{0.677}
\newcommand{\qwenCoCoAConf}{0.943}
\newcommand{\qwenCoCoACalGap}{0.266}
\newcommand{\qwenCoCoAECE}{0.266}
\newcommand{\qwenCoCoABrier}{0.267}
\newcommand{\qwenCoCoAHLCC}{0.801}
\newcommand{\qwenCoCoAAUROC}{0.806}
\newcommand{\qwenCoCoAAUPRC}{0.904}
\newcommand{\qwenCoCoAPRR}{0.726}
\newcommand{\qwenCoCoAAccSeventy}{0.810}
\newcommand{\qwenCoCoAAccEighty}{0.762}
\newcommand{\qwenCoCoAAccNinety}{0.715}

% Llama-3.1-8B (14 methods, 300 examples)
\newcommand{\llamaCalNExamples}{300}
\newcommand{\llamaCalAcc}{0.393}
\newcommand{\llamaEntropyAcc}{0.393}
\newcommand{\llamaEntropyConf}{0.917}
\newcommand{\llamaEntropyCalGap}{0.524}
\newcommand{\llamaEntropyECE}{0.524}
\newcommand{\llamaEntropyBrier}{0.491}
\newcommand{\llamaEntropyHLCC}{-0.211}
\newcommand{\llamaEntropyAUROC}{0.809}
\newcommand{\llamaEntropyAUPRC}{0.741}
\newcommand{\llamaEntropyPRR}{0.653}
\newcommand{\llamaEntropyAccSeventy}{0.529}
\newcommand{\llamaEntropyAccEighty}{0.471}
\newcommand{\llamaEntropyAccNinety}{0.430}
\newcommand{\llamaTopkAcc}{0.393}
\newcommand{\llamaTopkConf}{0.962}
\newcommand{\llamaTopkCalGap}{0.569}
\newcommand{\llamaTopkECE}{0.569}
\newcommand{\llamaTopkBrier}{0.548}
\newcommand{\llamaTopkHLCC}{-0.314}
\newcommand{\llamaTopkAUROC}{0.746}
\newcommand{\llamaTopkAUPRC}{0.626}
\newcommand{\llamaTopkPRR}{0.481}
\newcommand{\llamaTopkAccSeventy}{0.505}
\newcommand{\llamaTopkAccEighty}{0.458}
\newcommand{\llamaTopkAccNinety}{0.430}
\newcommand{\llamaMSPAcc}{0.393}
\newcommand{\llamaMSPConf}{0.683}
\newcommand{\llamaMSPCalGap}{0.290}
\newcommand{\llamaMSPECE}{0.295}
\newcommand{\llamaMSPBrier}{0.312}
\newcommand{\llamaMSPHLCC}{0.144}
\newcommand{\llamaMSPAUROC}{0.658}
\newcommand{\llamaMSPAUPRC}{0.527}
\newcommand{\llamaMSPPRR}{0.302}
\newcommand{\llamaMSPAccSeventy}{0.462}
\newcommand{\llamaMSPAccEighty}{0.438}
\newcommand{\llamaMSPAccNinety}{0.415}
\newcommand{\llamaGNLLAcc}{0.393}
\newcommand{\llamaGNLLConf}{0.683}
\newcommand{\llamaGNLLCalGap}{0.290}
\newcommand{\llamaGNLLECE}{0.295}
\newcommand{\llamaGNLLBrier}{0.312}
\newcommand{\llamaGNLLHLCC}{0.144}
\newcommand{\llamaGNLLAUROC}{0.658}
\newcommand{\llamaGNLLAUPRC}{0.527}
\newcommand{\llamaGNLLPRR}{0.302}
\newcommand{\llamaGNLLAccSeventy}{0.462}
\newcommand{\llamaGNLLAccEighty}{0.438}
\newcommand{\llamaGNLLAccNinety}{0.415}
\newcommand{\llamaVerbAcc}{0.393}
\newcommand{\llamaVerbConf}{0.667}
\newcommand{\llamaVerbCalGap}{0.274}
\newcommand{\llamaVerbECE}{0.421}
\newcommand{\llamaVerbBrier}{0.565}
\newcommand{\llamaVerbHLCC}{-0.198}
\newcommand{\llamaVerbAUROC}{0.467}
\newcommand{\llamaVerbAUPRC}{0.371}
\newcommand{\llamaVerbPRR}{-0.069}
\newcommand{\llamaVerbAccSeventy}{0.371}
\newcommand{\llamaVerbAccEighty}{0.383}
\newcommand{\llamaVerbAccNinety}{0.385}
\newcommand{\llamaSelfConAcc}{0.507}
\newcommand{\llamaSelfConConf}{0.748}
\newcommand{\llamaSelfConCalGap}{0.242}
\newcommand{\llamaSelfConECE}{0.242}
\newcommand{\llamaSelfConBrier}{0.279}
\newcommand{\llamaSelfConHLCC}{0.440}
\newcommand{\llamaSelfConAUROC}{0.707}
\newcommand{\llamaSelfConAUPRC}{0.716}
\newcommand{\llamaSelfConPRR}{0.485}
\newcommand{\llamaSelfConAccSeventy}{0.571}
\newcommand{\llamaSelfConAccEighty}{0.550}
\newcommand{\llamaSelfConAccNinety}{0.530}
\newcommand{\llamaTempScaleAcc}{0.393}
\newcommand{\llamaTempScaleConf}{0.960}
\newcommand{\llamaTempScaleCalGap}{0.566}
\newcommand{\llamaTempScaleECE}{0.566}
\newcommand{\llamaTempScaleBrier}{0.558}
\newcommand{\llamaTempScaleHLCC}{-0.333}
\newcommand{\llamaTempScaleAUROC}{0.636}
\newcommand{\llamaTempScaleAUPRC}{0.485}
\newcommand{\llamaTempScalePRR}{0.217}
\newcommand{\llamaTempScaleAccSeventy}{0.438}
\newcommand{\llamaTempScaleAccEighty}{0.429}
\newcommand{\llamaTempScaleAccNinety}{0.407}
\newcommand{\llamaIsotonicAcc}{0.393}
\newcommand{\llamaIsotonicConf}{0.393}
\newcommand{\llamaIsotonicCalGap}{0.000}
\newcommand{\llamaIsotonicECE}{0.000}
\newcommand{\llamaIsotonicBrier}{0.161}
\newcommand{\llamaIsotonicHLCC}{0.481}
\newcommand{\llamaIsotonicAUROC}{0.828}
\newcommand{\llamaIsotonicAUPRC}{0.751}
\newcommand{\llamaIsotonicPRR}{0.668}
\newcommand{\llamaIsotonicAccSeventy}{0.529}
\newcommand{\llamaIsotonicAccEighty}{0.479}
\newcommand{\llamaIsotonicAccNinety}{0.433}
\newcommand{\llamaPlattAcc}{0.393}
\newcommand{\llamaPlattConf}{0.393}
\newcommand{\llamaPlattCalGap}{0.000}
\newcommand{\llamaPlattECE}{0.075}
\newcommand{\llamaPlattBrier}{0.177}
\newcommand{\llamaPlattHLCC}{0.469}
\newcommand{\llamaPlattAUROC}{0.809}
\newcommand{\llamaPlattAUPRC}{0.741}
\newcommand{\llamaPlattPRR}{0.653}
\newcommand{\llamaPlattAccSeventy}{0.529}
\newcommand{\llamaPlattAccEighty}{0.471}
\newcommand{\llamaPlattAccNinety}{0.430}
\newcommand{\llamaNormShiftAcc}{0.393}
\newcommand{\llamaNormShiftConf}{0.534}
\newcommand{\llamaNormShiftCalGap}{0.141}
\newcommand{\llamaNormShiftECE}{0.235}
\newcommand{\llamaNormShiftBrier}{0.253}
\newcommand{\llamaNormShiftHLCC}{0.324}
\newcommand{\llamaNormShiftAUROC}{0.794}
\newcommand{\llamaNormShiftAUPRC}{0.713}
\newcommand{\llamaNormShiftPRR}{0.613}
\newcommand{\llamaNormShiftAccSeventy}{0.514}
\newcommand{\llamaNormShiftAccEighty}{0.471}
\newcommand{\llamaNormShiftAccNinety}{0.433}
\newcommand{\llamaSemEntAcc}{0.507}
\newcommand{\llamaSemEntConf}{0.650}
\newcommand{\llamaSemEntCalGap}{0.143}
\newcommand{\llamaSemEntECE}{0.205}
\newcommand{\llamaSemEntBrier}{0.253}
\newcommand{\llamaSemEntHLCC}{0.526}
\newcommand{\llamaSemEntAUROC}{0.691}
\newcommand{\llamaSemEntAUPRC}{0.744}
\newcommand{\llamaSemEntPRR}{0.527}
\newcommand{\llamaSemEntAccSeventy}{0.595}
\newcommand{\llamaSemEntAccEighty}{0.550}
\newcommand{\llamaSemEntAccNinety}{0.522}
\newcommand{\llamaCotConfAcc}{0.647}
\newcommand{\llamaCotConfConf}{0.807}
\newcommand{\llamaCotConfCalGap}{0.160}
\newcommand{\llamaCotConfECE}{0.213}
\newcommand{\llamaCotConfBrier}{0.261}
\newcommand{\llamaCotConfHLCC}{0.744}
\newcommand{\llamaCotConfAUROC}{0.630}
\newcommand{\llamaCotConfAUPRC}{0.741}
\newcommand{\llamaCotConfPRR}{0.293}
\newcommand{\llamaCotConfAccSeventy}{0.657}
\newcommand{\llamaCotConfAccEighty}{0.650}
\newcommand{\llamaCotConfAccNinety}{0.656}
\newcommand{\llamaLVUAcc}{0.393}
\newcommand{\llamaLVUConf}{0.453}
\newcommand{\llamaLVUCalGap}{0.060}
\newcommand{\llamaLVUECE}{0.120}
\newcommand{\llamaLVUBrier}{0.251}
\newcommand{\llamaLVUHLCC}{0.314}
\newcommand{\llamaLVUAUROC}{0.545}
\newcommand{\llamaLVUAUPRC}{0.459}
\newcommand{\llamaLVUPRR}{0.143}
\newcommand{\llamaLVUAccSeventy}{0.424}
\newcommand{\llamaLVUAccEighty}{0.412}
\newcommand{\llamaLVUAccNinety}{0.404}
\newcommand{\llamaCoCoAAcc}{0.523}
\newcommand{\llamaCoCoAConf}{0.825}
\newcommand{\llamaCoCoACalGap}{0.301}
\newcommand{\llamaCoCoAECE}{0.301}
\newcommand{\llamaCoCoABrier}{0.309}
\newcommand{\llamaCoCoAHLCC}{0.393}
\newcommand{\llamaCoCoAAUROC}{0.722}
\newcommand{\llamaCoCoAAUPRC}{0.777}
\newcommand{\llamaCoCoAPRR}{0.577}
\newcommand{\llamaCoCoAAccSeventy}{0.605}
\newcommand{\llamaCoCoAAccEighty}{0.583}
\newcommand{\llamaCoCoAAccNinety}{0.537}

% Qwen2.5-14B (14 methods, 300 examples)
\newcommand{\qwenLCalNExamples}{300}
\newcommand{\qwenLCalAcc}{0.620}
\newcommand{\qwenLEntropyAcc}{0.620}
\newcommand{\qwenLEntropyConf}{0.942}
\newcommand{\qwenLEntropyCalGap}{0.322}
\newcommand{\qwenLEntropyECE}{0.322}
\newcommand{\qwenLEntropyBrier}{0.311}
\newcommand{\qwenLEntropyHLCC}{0.600}
\newcommand{\qwenLEntropyAUROC}{0.901}
\newcommand{\qwenLEntropyAUPRC}{0.932}
\newcommand{\qwenLEntropyPRR}{0.845}
\newcommand{\qwenLEntropyAccSeventy}{0.805}
\newcommand{\qwenLEntropyAccEighty}{0.746}
\newcommand{\qwenLEntropyAccNinety}{0.689}
\newcommand{\qwenLTopkAcc}{0.620}
\newcommand{\qwenLTopkConf}{0.972}
\newcommand{\qwenLTopkCalGap}{0.352}
\newcommand{\qwenLTopkECE}{0.352}
\newcommand{\qwenLTopkBrier}{0.342}
\newcommand{\qwenLTopkHLCC}{0.549}
\newcommand{\qwenLTopkAUROC}{0.799}
\newcommand{\qwenLTopkAUPRC}{0.838}
\newcommand{\qwenLTopkPRR}{0.624}
\newcommand{\qwenLTopkAccSeventy}{0.786}
\newcommand{\qwenLTopkAccEighty}{0.742}
\newcommand{\qwenLTopkAccNinety}{0.681}
\newcommand{\qwenLMSPAcc}{0.620}
\newcommand{\qwenLMSPConf}{0.860}
\newcommand{\qwenLMSPCalGap}{0.240}
\newcommand{\qwenLMSPECE}{0.240}
\newcommand{\qwenLMSPBrier}{0.263}
\newcommand{\qwenLMSPHLCC}{0.694}
\newcommand{\qwenLMSPAUROC}{0.719}
\newcommand{\qwenLMSPAUPRC}{0.796}
\newcommand{\qwenLMSPPRR}{0.508}
\newcommand{\qwenLMSPAccSeventy}{0.733}
\newcommand{\qwenLMSPAccEighty}{0.688}
\newcommand{\qwenLMSPAccNinety}{0.652}
\newcommand{\qwenLGNLLAcc}{0.620}
\newcommand{\qwenLGNLLConf}{0.860}
\newcommand{\qwenLGNLLCalGap}{0.240}
\newcommand{\qwenLGNLLECE}{0.240}
\newcommand{\qwenLGNLLBrier}{0.263}
\newcommand{\qwenLGNLLHLCC}{0.694}
\newcommand{\qwenLGNLLAUROC}{0.719}
\newcommand{\qwenLGNLLAUPRC}{0.796}
\newcommand{\qwenLGNLLPRR}{0.508}
\newcommand{\qwenLGNLLAccSeventy}{0.733}
\newcommand{\qwenLGNLLAccEighty}{0.688}
\newcommand{\qwenLGNLLAccNinety}{0.652}
\newcommand{\qwenLVerbAcc}{0.620}
\newcommand{\qwenLVerbConf}{0.570}
\newcommand{\qwenLVerbCalGap}{0.050}
\newcommand{\qwenLVerbECE}{0.229}
\newcommand{\qwenLVerbBrier}{0.298}
\newcommand{\qwenLVerbHLCC}{0.731}
\newcommand{\qwenLVerbAUROC}{0.604}
\newcommand{\qwenLVerbAUPRC}{0.707}
\newcommand{\qwenLVerbPRR}{0.264}
\newcommand{\qwenLVerbAccSeventy}{0.667}
\newcommand{\qwenLVerbAccEighty}{0.642}
\newcommand{\qwenLVerbAccNinety}{0.637}
\newcommand{\qwenLSelfConAcc}{0.740}
\newcommand{\qwenLSelfConConf}{0.916}
\newcommand{\qwenLSelfConCalGap}{0.176}
\newcommand{\qwenLSelfConECE}{0.178}
\newcommand{\qwenLSelfConBrier}{0.180}
\newcommand{\qwenLSelfConHLCC}{1.113}
\newcommand{\qwenLSelfConAUROC}{0.747}
\newcommand{\qwenLSelfConAUPRC}{0.859}
\newcommand{\qwenLSelfConPRR}{0.490}
\newcommand{\qwenLSelfConAccSeventy}{0.862}
\newcommand{\qwenLSelfConAccEighty}{0.846}
\newcommand{\qwenLSelfConAccNinety}{0.796}
\newcommand{\qwenLTempScaleAcc}{0.620}
\newcommand{\qwenLTempScaleConf}{0.986}
\newcommand{\qwenLTempScaleCalGap}{0.366}
\newcommand{\qwenLTempScaleECE}{0.366}
\newcommand{\qwenLTempScaleBrier}{0.361}
\newcommand{\qwenLTempScaleHLCC}{0.516}
\newcommand{\qwenLTempScaleAUROC}{0.669}
\newcommand{\qwenLTempScaleAUPRC}{0.762}
\newcommand{\qwenLTempScalePRR}{0.420}
\newcommand{\qwenLTempScaleAccSeventy}{0.733}
\newcommand{\qwenLTempScaleAccEighty}{0.675}
\newcommand{\qwenLTempScaleAccNinety}{0.641}
\newcommand{\qwenLIsotonicAcc}{0.620}
\newcommand{\qwenLIsotonicConf}{0.620}
\newcommand{\qwenLIsotonicCalGap}{0.000}
\newcommand{\qwenLIsotonicECE}{0.000}
\newcommand{\qwenLIsotonicBrier}{0.112}
\newcommand{\qwenLIsotonicHLCC}{1.007}
\newcommand{\qwenLIsotonicAUROC}{0.910}
\newcommand{\qwenLIsotonicAUPRC}{0.941}
\newcommand{\qwenLIsotonicPRR}{0.867}
\newcommand{\qwenLIsotonicAccSeventy}{0.810}
\newcommand{\qwenLIsotonicAccEighty}{0.750}
\newcommand{\qwenLIsotonicAccNinety}{0.689}
\newcommand{\qwenLPlattAcc}{0.620}
\newcommand{\qwenLPlattConf}{0.620}
\newcommand{\qwenLPlattCalGap}{0.000}
\newcommand{\qwenLPlattECE}{0.057}
\newcommand{\qwenLPlattBrier}{0.122}
\newcommand{\qwenLPlattHLCC}{0.984}
\newcommand{\qwenLPlattAUROC}{0.901}
\newcommand{\qwenLPlattAUPRC}{0.932}
\newcommand{\qwenLPlattPRR}{0.845}
\newcommand{\qwenLPlattAccSeventy}{0.805}
\newcommand{\qwenLPlattAccEighty}{0.746}
\newcommand{\qwenLPlattAccNinety}{0.689}
\newcommand{\qwenLNormShiftAcc}{0.620}
\newcommand{\qwenLNormShiftConf}{0.786}
\newcommand{\qwenLNormShiftCalGap}{0.166}
\newcommand{\qwenLNormShiftECE}{0.183}
\newcommand{\qwenLNormShiftBrier}{0.201}
\newcommand{\qwenLNormShiftHLCC}{0.840}
\newcommand{\qwenLNormShiftAUROC}{0.786}
\newcommand{\qwenLNormShiftAUPRC}{0.811}
\newcommand{\qwenLNormShiftPRR}{0.558}
\newcommand{\qwenLNormShiftAccSeventy}{0.786}
\newcommand{\qwenLNormShiftAccEighty}{0.733}
\newcommand{\qwenLNormShiftAccNinety}{0.674}
\newcommand{\qwenLSemEntAcc}{0.747}
\newcommand{\qwenLSemEntConf}{0.889}
\newcommand{\qwenLSemEntCalGap}{0.142}
\newcommand{\qwenLSemEntECE}{0.171}
\newcommand{\qwenLSemEntBrier}{0.177}
\newcommand{\qwenLSemEntHLCC}{1.140}
\newcommand{\qwenLSemEntAUROC}{0.740}
\newcommand{\qwenLSemEntAUPRC}{0.876}
\newcommand{\qwenLSemEntPRR}{0.542}
\newcommand{\qwenLSemEntAccSeventy}{0.871}
\newcommand{\qwenLSemEntAccEighty}{0.829}
\newcommand{\qwenLSemEntAccNinety}{0.778}
\newcommand{\qwenLCotConfAcc}{0.767}
\newcommand{\qwenLCotConfConf}{0.879}
\newcommand{\qwenLCotConfCalGap}{0.112}
\newcommand{\qwenLCotConfECE}{0.189}
\newcommand{\qwenLCotConfBrier}{0.197}
\newcommand{\qwenLCotConfHLCC}{1.114}
\newcommand{\qwenLCotConfAUROC}{0.689}
\newcommand{\qwenLCotConfAUPRC}{0.875}
\newcommand{\qwenLCotConfPRR}{0.486}
\newcommand{\qwenLCotConfAccSeventy}{0.848}
\newcommand{\qwenLCotConfAccEighty}{0.792}
\newcommand{\qwenLCotConfAccNinety}{0.767}
\newcommand{\qwenLLVUAcc}{0.620}
\newcommand{\qwenLLVUConf}{0.358}
\newcommand{\qwenLLVUCalGap}{0.262}
\newcommand{\qwenLLVUECE}{0.262}
\newcommand{\qwenLLVUBrier}{0.328}
\newcommand{\qwenLLVUHLCC}{0.736}
\newcommand{\qwenLLVUAUROC}{0.531}
\newcommand{\qwenLLVUAUPRC}{0.667}
\newcommand{\qwenLLVUPRR}{0.137}
\newcommand{\qwenLLVUAccSeventy}{0.629}
\newcommand{\qwenLLVUAccEighty}{0.625}
\newcommand{\qwenLLVUAccNinety}{0.619}
\newcommand{\qwenLCoCoAAcc}{0.740}
\newcommand{\qwenLCoCoAConf}{0.923}
\newcommand{\qwenLCoCoACalGap}{0.183}
\newcommand{\qwenLCoCoAECE}{0.183}
\newcommand{\qwenLCoCoABrier}{0.188}
\newcommand{\qwenLCoCoAHLCC}{1.082}
\newcommand{\qwenLCoCoAAUROC}{0.846}
\newcommand{\qwenLCoCoAAUPRC}{0.936}
\newcommand{\qwenLCoCoAPRR}{0.773}
\newcommand{\qwenLCoCoAAccSeventy}{0.881}
\newcommand{\qwenLCoCoAAccEighty}{0.846}
\newcommand{\qwenLCoCoAAccNinety}{0.793}

% DeepSeek-R1-14B (14 methods, 300 examples)
\newcommand{\deepseekCalNExamples}{300}
\newcommand{\deepseekCalAcc}{0.497}
\newcommand{\deepseekEntropyAcc}{0.497}
\newcommand{\deepseekEntropyConf}{0.869}
\newcommand{\deepseekEntropyCalGap}{0.373}
\newcommand{\deepseekEntropyECE}{0.373}
\newcommand{\deepseekEntropyBrier}{0.364}
\newcommand{\deepseekEntropyHLCC}{0.225}
\newcommand{\deepseekEntropyAUROC}{0.839}
\newcommand{\deepseekEntropyAUPRC}{0.838}
\newcommand{\deepseekEntropyPRR}{0.731}
\newcommand{\deepseekEntropyAccSeventy}{0.633}
\newcommand{\deepseekEntropyAccEighty}{0.596}
\newcommand{\deepseekEntropyAccNinety}{0.552}
\newcommand{\deepseekTopkAcc}{0.497}
\newcommand{\deepseekTopkConf}{0.880}
\newcommand{\deepseekTopkCalGap}{0.383}
\newcommand{\deepseekTopkECE}{0.383}
\newcommand{\deepseekTopkBrier}{0.374}
\newcommand{\deepseekTopkHLCC}{0.211}
\newcommand{\deepseekTopkAUROC}{0.742}
\newcommand{\deepseekTopkAUPRC}{0.748}
\newcommand{\deepseekTopkPRR}{0.560}
\newcommand{\deepseekTopkAccSeventy}{0.600}
\newcommand{\deepseekTopkAccEighty}{0.575}
\newcommand{\deepseekTopkAccNinety}{0.533}
\newcommand{\deepseekMSPAcc}{0.497}
\newcommand{\deepseekMSPConf}{0.714}
\newcommand{\deepseekMSPCalGap}{0.218}
\newcommand{\deepseekMSPECE}{0.234}
\newcommand{\deepseekMSPBrier}{0.302}
\newcommand{\deepseekMSPHLCC}{0.365}
\newcommand{\deepseekMSPAUROC}{0.602}
\newcommand{\deepseekMSPAUPRC}{0.552}
\newcommand{\deepseekMSPPRR}{0.143}
\newcommand{\deepseekMSPAccSeventy}{0.562}
\newcommand{\deepseekMSPAccEighty}{0.525}
\newcommand{\deepseekMSPAccNinety}{0.507}
\newcommand{\deepseekGNLLAcc}{0.497}
\newcommand{\deepseekGNLLConf}{0.714}
\newcommand{\deepseekGNLLCalGap}{0.218}
\newcommand{\deepseekGNLLECE}{0.234}
\newcommand{\deepseekGNLLBrier}{0.302}
\newcommand{\deepseekGNLLHLCC}{0.365}
\newcommand{\deepseekGNLLAUROC}{0.602}
\newcommand{\deepseekGNLLAUPRC}{0.552}
\newcommand{\deepseekGNLLPRR}{0.143}
\newcommand{\deepseekGNLLAccSeventy}{0.562}
\newcommand{\deepseekGNLLAccEighty}{0.525}
\newcommand{\deepseekGNLLAccNinety}{0.507}
\newcommand{\deepseekVerbAcc}{0.497}
\newcommand{\deepseekVerbConf}{0.761}
\newcommand{\deepseekVerbCalGap}{0.264}
\newcommand{\deepseekVerbECE}{0.347}
\newcommand{\deepseekVerbBrier}{0.371}
\newcommand{\deepseekVerbHLCC}{0.256}
\newcommand{\deepseekVerbAUROC}{0.560}
\newcommand{\deepseekVerbAUPRC}{0.542}
\newcommand{\deepseekVerbPRR}{0.121}
\newcommand{\deepseekVerbAccSeventy}{0.524}
\newcommand{\deepseekVerbAccEighty}{0.517}
\newcommand{\deepseekVerbAccNinety}{0.496}
\newcommand{\deepseekSelfConAcc}{0.983}
\newcommand{\deepseekSelfConConf}{0.973}
\newcommand{\deepseekSelfConCalGap}{0.010}
\newcommand{\deepseekSelfConECE}{0.013}
\newcommand{\deepseekSelfConBrier}{0.012}
\newcommand{\deepseekSelfConHLCC}{1.933}
\newcommand{\deepseekSelfConAUROC}{0.984}
\newcommand{\deepseekSelfConAUPRC}{1.000}
\newcommand{\deepseekSelfConPRR}{0.985}
\newcommand{\deepseekSelfConAccSeventy}{1.000}
\newcommand{\deepseekSelfConAccEighty}{1.000}
\newcommand{\deepseekSelfConAccNinety}{1.000}
\newcommand{\deepseekTempScaleAcc}{0.497}
\newcommand{\deepseekTempScaleConf}{0.970}
\newcommand{\deepseekTempScaleCalGap}{0.473}
\newcommand{\deepseekTempScaleECE}{0.477}
\newcommand{\deepseekTempScaleBrier}{0.480}
\newcommand{\deepseekTempScaleHLCC}{0.029}
\newcommand{\deepseekTempScaleAUROC}{0.597}
\newcommand{\deepseekTempScaleAUPRC}{0.603}
\newcommand{\deepseekTempScalePRR}{0.262}
\newcommand{\deepseekTempScaleAccSeventy}{0.562}
\newcommand{\deepseekTempScaleAccEighty}{0.529}
\newcommand{\deepseekTempScaleAccNinety}{0.500}
\newcommand{\deepseekIsotonicAcc}{0.497}
\newcommand{\deepseekIsotonicConf}{0.497}
\newcommand{\deepseekIsotonicCalGap}{0.000}
\newcommand{\deepseekIsotonicECE}{0.000}
\newcommand{\deepseekIsotonicBrier}{0.153}
\newcommand{\deepseekIsotonicHLCC}{0.691}
\newcommand{\deepseekIsotonicAUROC}{0.853}
\newcommand{\deepseekIsotonicAUPRC}{0.856}
\newcommand{\deepseekIsotonicPRR}{0.760}
\newcommand{\deepseekIsotonicAccSeventy}{0.638}
\newcommand{\deepseekIsotonicAccEighty}{0.600}
\newcommand{\deepseekIsotonicAccNinety}{0.552}
\newcommand{\deepseekPlattAcc}{0.497}
\newcommand{\deepseekPlattConf}{0.497}
\newcommand{\deepseekPlattCalGap}{0.000}
\newcommand{\deepseekPlattECE}{0.045}
\newcommand{\deepseekPlattBrier}{0.164}
\newcommand{\deepseekPlattHLCC}{0.665}
\newcommand{\deepseekPlattAUROC}{0.839}
\newcommand{\deepseekPlattAUPRC}{0.838}
\newcommand{\deepseekPlattPRR}{0.731}
\newcommand{\deepseekPlattAccSeventy}{0.633}
\newcommand{\deepseekPlattAccEighty}{0.596}
\newcommand{\deepseekPlattAccNinety}{0.552}
\newcommand{\deepseekNormShiftAcc}{0.497}
\newcommand{\deepseekNormShiftConf}{0.479}
\newcommand{\deepseekNormShiftCalGap}{0.018}
\newcommand{\deepseekNormShiftECE}{0.112}
\newcommand{\deepseekNormShiftBrier}{0.113}
\newcommand{\deepseekNormShiftHLCC}{0.819}
\newcommand{\deepseekNormShiftAUROC}{0.948}
\newcommand{\deepseekNormShiftAUPRC}{0.933}
\newcommand{\deepseekNormShiftPRR}{0.898}
\newcommand{\deepseekNormShiftAccSeventy}{0.710}
\newcommand{\deepseekNormShiftAccEighty}{0.621}
\newcommand{\deepseekNormShiftAccNinety}{0.552}
\newcommand{\deepseekSemEntAcc}{0.990}
\newcommand{\deepseekSemEntConf}{0.962}
\newcommand{\deepseekSemEntCalGap}{0.028}
\newcommand{\deepseekSemEntECE}{0.028}
\newcommand{\deepseekSemEntBrier}{0.017}
\newcommand{\deepseekSemEntHLCC}{1.937}
\newcommand{\deepseekSemEntAUROC}{0.961}
\newcommand{\deepseekSemEntAUPRC}{1.000}
\newcommand{\deepseekSemEntPRR}{0.960}
\newcommand{\deepseekSemEntAccSeventy}{1.000}
\newcommand{\deepseekSemEntAccEighty}{1.000}
\newcommand{\deepseekSemEntAccNinety}{1.000}
\newcommand{\deepseekCotConfAcc}{1.000}
\newcommand{\deepseekCotConfConf}{0.500}
\newcommand{\deepseekCotConfCalGap}{0.500}
\newcommand{\deepseekCotConfECE}{0.500}
\newcommand{\deepseekCotConfBrier}{0.250}
\newcommand{\deepseekCotConfHLCC}{1.500}
\newcommand{\deepseekCotConfAUROC}{0.500}
\newcommand{\deepseekCotConfAUPRC}{0.500}
\newcommand{\deepseekCotConfPRR}{0.000}
\newcommand{\deepseekCotConfAccSeventy}{1.000}
\newcommand{\deepseekCotConfAccEighty}{1.000}
\newcommand{\deepseekCotConfAccNinety}{1.000}
\newcommand{\deepseekLVUAcc}{0.497}
\newcommand{\deepseekLVUConf}{0.251}
\newcommand{\deepseekLVUCalGap}{0.246}
\newcommand{\deepseekLVUECE}{0.256}
\newcommand{\deepseekLVUBrier}{0.351}
\newcommand{\deepseekLVUHLCC}{0.517}
\newcommand{\deepseekLVUAUROC}{0.484}
\newcommand{\deepseekLVUAUPRC}{0.504}
\newcommand{\deepseekLVUPRR}{0.007}
\newcommand{\deepseekLVUAccSeventy}{0.486}
\newcommand{\deepseekLVUAccEighty}{0.500}
\newcommand{\deepseekLVUAccNinety}{0.496}
\newcommand{\deepseekCoCoAAcc}{0.987}
\newcommand{\deepseekCoCoAConf}{0.923}
\newcommand{\deepseekCoCoACalGap}{0.064}
\newcommand{\deepseekCoCoAECE}{0.071}
\newcommand{\deepseekCoCoABrier}{0.015}
\newcommand{\deepseekCoCoAHLCC}{1.884}
\newcommand{\deepseekCoCoAAUROC}{0.976}
\newcommand{\deepseekCoCoAAUPRC}{1.000}
\newcommand{\deepseekCoCoAPRR}{0.975}
\newcommand{\deepseekCoCoAAccSeventy}{1.000}
\newcommand{\deepseekCoCoAAccEighty}{1.000}
\newcommand{\deepseekCoCoAAccNinety}{1.000}

% ============================================
% Per-Layer Signal Analysis
% SOURCE: D:/git/NNConfidence\data/results/<model>_layer_signals.json
% ============================================
% Mistral-7B: 32 layers, 90 examples
\newcommand{\mistralNLayers}{32}
\newcommand{\mistralHiddenSize}{4096}
\newcommand{\mistralLayerAcc}{66.7\%}
\newcommand{\mistralBestLayer}{27}
\newcommand{\mistralBestLayerD}{0.82}
\newcommand{\mistralSecondLayer}{25}
\newcommand{\mistralSecondLayerD}{0.72}
\newcommand{\mistralThirdLayer}{24}
\newcommand{\mistralThirdLayerD}{0.68}

% Qwen2.5-7B: 28 layers, 90 examples
\newcommand{\qwenNLayers}{28}
\newcommand{\qwenHiddenSize}{3584}
\newcommand{\qwenLayerAcc}{86.7\%}
\newcommand{\qwenBestLayer}{22}
\newcommand{\qwenBestLayerD}{0.76}
\newcommand{\qwenSecondLayer}{27}
\newcommand{\qwenSecondLayerD}{0.74}
\newcommand{\qwenThirdLayer}{15}
\newcommand{\qwenThirdLayerD}{-0.67}

% Llama-3.1-8B: 32 layers, 90 examples
\newcommand{\llamaNLayers}{32}
\newcommand{\llamaHiddenSize}{4096}
\newcommand{\llamaLayerAcc}{72.2\%}
\newcommand{\llamaBestLayer}{2}
\newcommand{\llamaBestLayerD}{0.99}
\newcommand{\llamaSecondLayer}{3}
\newcommand{\llamaSecondLayerD}{-0.90}
\newcommand{\llamaThirdLayer}{14}
\newcommand{\llamaThirdLayerD}{-0.90}

% Qwen2.5-14B: 48 layers, 90 examples
\newcommand{\qwenLNLayers}{48}
\newcommand{\qwenLHiddenSize}{5120}
\newcommand{\qwenLLayerAcc}{92.2\%}
\newcommand{\qwenLBestLayer}{30}
\newcommand{\qwenLBestLayerD}{-1.53}
\newcommand{\qwenLSecondLayer}{47}
\newcommand{\qwenLSecondLayerD}{1.24}
\newcommand{\qwenLThirdLayer}{31}
\newcommand{\qwenLThirdLayerD}{-1.11}

% DeepSeek-R1-14B: 48 layers, 90 examples
\newcommand{\deepseekNLayers}{48}
\newcommand{\deepseekHiddenSize}{5120}
\newcommand{\deepseekLayerAcc}{66.7\%}
\newcommand{\deepseekBestLayer}{15}
\newcommand{\deepseekBestLayerD}{-0.85}
\newcommand{\deepseekSecondLayer}{27}
\newcommand{\deepseekSecondLayerD}{-0.83}
\newcommand{\deepseekThirdLayer}{9}
\newcommand{\deepseekThirdLayerD}{0.77}

% ============================================
% Model Metadata
% SOURCE: D:/git/NNConfidence\paper_constants.yaml
% ============================================
\newcommand{\mistralDisplayName}{Mistral-7B}
\newcommand{\mistralParams}{7.2B}
\newcommand{\mistralPrecision}{FP16}
\newcommand{\qwenDisplayName}{Qwen2.5-7B}
\newcommand{\qwenParams}{7.6B}
\newcommand{\qwenPrecision}{FP16}
\newcommand{\llamaDisplayName}{Llama-3.1-8B}
\newcommand{\llamaParams}{8.0B}
\newcommand{\llamaPrecision}{FP16}
\newcommand{\qwenLDisplayName}{Qwen2.5-14B}
\newcommand{\qwenLParams}{14.2B}
\newcommand{\qwenLPrecision}{8-bit}
\newcommand{\deepseekDisplayName}{DeepSeek-R1-14B}
\newcommand{\deepseekParams}{14.2B}
\newcommand{\deepseekPrecision}{8-bit}
\newcommand{\tinyllamaDisplayName}{TinyLlama}
\newcommand{\tinyllamaParams}{1.1B}
\newcommand{\tinyllamaPrecision}{FP16}
\newcommand{\qwenSmallDisplayName}{Qwen2-1.5B}
\newcommand{\qwenSmallParams}{1.5B}
\newcommand{\qwenSmallPrecision}{FP16}
\newcommand{\phiTwoDisplayName}{Phi-2}
\newcommand{\phiTwoParams}{2.7B}
\newcommand{\phiTwoPrecision}{FP16}
\newcommand{\phiThreeMiniDisplayName}{Phi-3-mini}
\newcommand{\phiThreeMiniParams}{3.8B}
\newcommand{\phiThreeMiniPrecision}{FP16}



\begin{abstract}
We introduce the Hybrid Linear-Convex Confidence (\hlcc{}) loss, a
scoring function that combines a linear reward for correct predictions
with a quadratic penalty for confident errors.  Unlike proper scoring
rules such as the Brier score or log-loss, where the optimal confidence
equals the true accuracy ($\cstar = p$), \hlcc{} induces a conservative
mapping $\cstar = p/[4(1-p)]$ that naturally suppresses overconfidence
below the crossover point at $p = 0.75$.  We implement \hlcc{} as a trainable objective for
language-model confidence heads, proposing a dual-path
\emph{NormShift} architecture that routes per-layer normalisation
dynamics alongside final hidden states to estimate per-token confidence.
We benchmark against 13 baseline confidence methods (including maximum
softmax probability, entropy, self-consistency, semantic entropy,
chain-of-thought verbalisation, temperature scaling, isotonic
regression, and Platt scaling) across 9 datasets and 9 model scales
from 1.1B to 14.2B parameters.  Results are evaluated on discrimination
(AUROC, AUPRC), calibration (ECE, Brier), and selective prediction
(PRR, accuracy at coverage thresholds), providing a comprehensive
assessment of confidence estimation in modern language models.
\end{abstract}

% ============================================================
\section{Introduction}
\label{sec:introduction}
% ============================================================

Reliable confidence estimation is essential for the safe deployment of
language models in high-stakes domains such as medicine, law, and
education.  Standard autoregressive training with cross-entropy loss
maximises token-level likelihood but provides no explicit training signal
for calibrated uncertainty~\cite{guo2017calibration}.  The resulting
softmax probabilities are often poorly calibrated, with models exhibiting
systematic overconfidence on out-of-distribution
inputs~\cite{hendrycks2017baseline}.

Proper scoring rules, notably the Brier score and logarithmic loss, are the
classical tools for eliciting calibrated probabilities.  Both are
\emph{strictly proper}: the expected score is maximised only when the
reported confidence equals the true probability of correctness.  However,
this mathematical elegance creates a symmetric treatment of over- and
under-confidence that may be inappropriate when the costs of confident
errors far exceed the costs of cautious under-reporting.

In certainty-based marking (CBM) for educational
assessment~\cite{gardnermedwin1995,gardnermedwin2006}, the asymmetry
between ``confident and wrong'' versus ``cautious and right'' is encoded
through discrete penalty scales derived from logarithmic proper scoring
rules~\cite{good1952}.  Inspired by this asymmetry, we propose the
Hybrid Linear-Convex Confidence (\hlcc{}) scoring function:
\begin{equation}
  S(y, c) =
  \begin{cases}
    1 + c    & \text{if } y = 1 \text{ (correct)}, \\
    -2c^{2}  & \text{if } y = 0 \text{ (incorrect)},
  \end{cases}
  \label{eq:hlcc-scoring}
\end{equation}
where $c \in [0,1]$ is the stated confidence and $y \in \{0,1\}$
indicates correctness.  The linear reward provides a constant marginal
incentive for increased confidence on correct predictions, while the
quadratic penalty creates an escalating marginal cost for confident
errors.  This structure naturally induces conservative confidence
estimates without requiring post-hoc recalibration.

\paragraph{Contributions.}
\begin{enumerate}[nosep]
  \item We derive the theoretical properties of \hlcc{}, showing that
        the optimal confidence $\cstar = p/[4(1-p)]$ is strictly below
        the true accuracy $p$ for $p < 0.75$, providing a built-in
        conservative bias (\Cref{sec:theory}).
  \item We propose the \emph{NormShift} confidence head, a dual-path
        architecture that combines per-layer normalisation shift signals
        with final hidden states to estimate token-level confidence
        (\Cref{sec:architecture}).
  \item We conduct a comprehensive benchmarking study covering 11
        baselines, 9 datasets, and 9 models from 1.1B to 14.2B
        parameters, with metrics spanning discrimination, calibration,
        and selective prediction
        (\Cref{sec:methodology,sec:results}).
\end{enumerate}


% ============================================================
\section{Theoretical Framework}
\label{sec:theory}
% ============================================================

\subsection{The HLCC Scoring Function}

\begin{definition}[HLCC Score]
  For a binary outcome $y \in \{0,1\}$ and stated confidence
  $c \in [0,1]$, the \hlcc{} score is
  \begin{equation}
    S(y, c) = y(1 + c) + (1-y)(-2c^{2}).
  \end{equation}
\end{definition}

The reward function $R(c) = 1 + c$ is affine with constant derivative
$R'(c) = 1$.  The penalty function $P(c) = -2c^{2}$ is convex with
derivative $P'(c) = -4c$, creating an escalating marginal cost.  At
$c = 0$, the score is $S(1,0) = 1$ and $S(0,0) = 0$, establishing a
\emph{safe guess} baseline: answering at zero confidence is always
weakly dominant over abstention.

\subsection{Optimal Confidence}

\begin{proposition}[Optimal Confidence under HLCC]
\label{prop:optimal}
  For an agent with true probability of correctness $p \in [0,1]$, the
  confidence $\cstar$ maximising expected \hlcc{} score is
  \begin{equation}
    \cstar = \min\!\left(1,\;\frac{p}{4(1-p)}\right).
    \label{eq:cstar}
  \end{equation}
\end{proposition}

\begin{proof}
  The expected score as a function of stated confidence $c$ is
  \begin{equation}
    \E[S] = p(1 + c) + (1-p)(-2c^{2}) = p + pc - 2(1-p)c^{2}.
  \end{equation}
  Differentiating with respect to $c$:
  \begin{equation}
    \frac{\partial \E[S]}{\partial c} = p - 4(1-p)c.
  \end{equation}
  Setting to zero and solving for $c$:
  \begin{equation}
    \cstar = \frac{p}{4(1-p)},
  \end{equation}
  subject to the constraint $c \in [0,1]$, giving the minimum with~$1$.
  The second derivative $\partial^{2}\E[S]/\partial c^{2} = -4(1-p) < 0$
  for $p < 1$, confirming this is a maximum.
\end{proof}

\begin{corollary}[Crossover Point]
\label{cor:crossover}
  The \hlcc{} optimal confidence equals the true accuracy ($\cstar = p$)
  at exactly $p = 0.75$.  For $p < 0.75$, $\cstar < p$ (conservative).
  For $p > 0.75$, $\cstar > p$ until saturation at $\cstar = 1$ when
  $p \geq 0.8$.
\end{corollary}

\begin{proof}
  Setting $\cstar = p$: $p = p/[4(1-p)]$ implies $4(1-p) = 1$, hence
  $p = 3/4$.  The saturation point solves $p/[4(1-p)] = 1$, giving
  $p = 4/5$.
\end{proof}

\begin{proposition}[Theoretical Maximum HLCC]
\label{prop:max-hlcc}
  For a model with accuracy $p$, the maximum achievable HLCC score is
  $S_{\max} = 2p$, attained when $c = 1$ for every correct prediction
  and $c = 0$ for every incorrect prediction.  The \emph{HLCC efficiency}
  $\eta = S_{\text{HLCC}} / (2p)$ measures how well the confidence method
  exploits the model's accuracy, with $\eta = 1$ indicating perfect
  self-knowledge.
\end{proposition}

\begin{proof}
  With perfect self-knowledge, each correct prediction ($p$ fraction)
  scores $S(1, 1) = 1 + 1 = 2$, and each incorrect prediction
  ($1 - p$ fraction) scores $S(0, 0) = 0$.  The expected score is
  $p \cdot 2 + (1-p) \cdot 0 = 2p$.  No other confidence assignment
  achieves a higher expected score, since $R(c) \leq 2$ and $P(c) \leq 0$.
\end{proof}

\Cref{tab:rational-confidence} shows the optimal confidence mapping at
representative accuracy levels.

\begin{table}[t]
\centering
\caption{Optimal confidence $\cstar$ under \hlcc{} compared with Brier/log-loss
         ($\cstar = p$) and CBM discrete levels.
         Data: \texttt{tables/rational\_confidence.csv}.}
\label{tab:rational-confidence}
\begin{tabular}{@{}lcccc@{}}
\toprule
True Accuracy $p$ & \hlcc{} $\cstar$ & Brier $\cstar$ & CBM Level & Behavioural Regime \\
\midrule
\csvreader[
  head to column names,
  late after line=\\,
  late after last line=\\,
]{tables/rational_confidence.csv}{}{%
  \p & \hlcccstar & \briercstar & \cbmlevel & \regime
}%
\bottomrule
\end{tabular}
\end{table}

\subsection{Comparison with Proper Scoring Rules}

\hlcc{} is \emph{not} a strictly proper scoring rule.  Unlike the Brier
score $S_B = -(c-y)^{2}$ and log-loss
$S_L = y\log c + (1-y)\log(1-c)$, which are both maximised at $\cstar = p$,
\hlcc{} imposes a conservative transformation of accuracy to confidence.

\begin{proposition}[Reward--Penalty Asymmetry]
  At any confidence level $c > 0$, the marginal penalty for an incorrect
  prediction exceeds the marginal reward for a correct one:
  \begin{equation}
    |P'(c)| = 4c > 1 = R'(c) \quad \text{for all } c > 0.25.
  \end{equation}
  The marginal risk-to-reward ratio is $4c : 1$.
\end{proposition}

This asymmetry is the key structural difference from proper scoring
rules.  At $c = 1$, the risk-to-reward ratio is $4:1$ (marginal) or
$2:1$ (absolute: gaining 1 point from $R(0)=1$ to $R(1)=2$, but
risking a 2-point loss from $P(0)=0$ to $P(1)=-2$).

\subsection{Safe Guess Mechanism}

At $c = 0$, the score is $S(1,0) = 1$ (correct) and $S(0,0) = 0$
(incorrect).  The expected value of a random guess ($p = 0.25$ for
4-option MCQ) at $c = 0$ is $0.25 \times 1 = 0.25 > 0$.  This ensures
that participating is always dominant over abstention, resolving the
\emph{omission bias} observed in negative-marking schemes where
risk-averse test-takers skip questions they partially
know~\cite{budescu2012optimal}.


% ============================================================
\section{Architecture}
\label{sec:architecture}
% ============================================================

\subsection{NormShift Confidence Head}

We propose the \emph{NormShift} confidence head, a lightweight module
trained on top of a frozen language model to predict per-token
confidence.  The architecture exploits an observation about transformer
normalisation layers: the magnitude of the correction applied by
LayerNorm encodes information about the model's internal uncertainty.

\begin{definition}[Norm-Shift Signal]
  For layer $\ell$ of a transformer with LayerNorm producing output
  $\mathbf{h}^{(\ell)}$ from pre-normalisation input
  $\mathbf{x}^{(\ell)}$, the norm-shift signal is
  \begin{equation}
    \delta^{(\ell)} = \|\mathbf{h}^{(\ell)}\|_{2} - \|\mathbf{x}^{(\ell)}\|_{2},
  \end{equation}
  measuring how much the normalisation operation altered the
  representation magnitude.
\end{definition}

\paragraph{Dual-Path Architecture.}
The NormShift head processes two complementary information streams:

\begin{itemize}[nosep]
  \item \textbf{Path A} (Norm-Shift Path): The vector
        $\boldsymbol{\delta} = [\delta^{(1)}, \ldots, \delta^{(L)}] \in \R^{L}$
        of per-layer norm-shift signals is projected through a linear
        layer with GELU activation to an intermediate representation of
        dimension $d_{\text{mid}} = \max(64, 2L)$.
  \item \textbf{Path B} (Hidden-State Path): The final hidden state
        $\mathbf{h}^{(L)} \in \R^{d}$ is projected through a separate
        linear layer with GELU activation to the same intermediate
        dimension $d_{\text{mid}}$.
\end{itemize}

The two paths are concatenated and processed by a combination network:
\begin{equation}
  c = \sigma\!\left(\frac{1}{\tau}\,\mathbf{w}^{\top}
  \text{GELU}\!\left(\mathbf{W}_{\text{comb}}
  \begin{bmatrix} f_{A}(\boldsymbol{\delta}) \\ f_{B}(\mathbf{h}^{(L)}) \end{bmatrix}
  + \mathbf{b}_{\text{comb}}\right)\right),
\end{equation}
where $\tau$ is a learnable temperature parameter initialised to 1, and
$\sigma$ is the sigmoid function.  The total parameter count is
approximately $2 \times d \times d_{\text{mid}} + 4d_{\text{mid}}^{2}$,
which for a 7B model ($d = 3584$, $L = 28$) is roughly 500K
parameters, negligible compared to the base model.

\paragraph{Training Objective.}
The confidence head is trained with the \hlcc{} loss from
\Cref{eq:hlcc-scoring}, combined with a cross-entropy warmup:
\begin{equation}
  \mathcal{L} = \alpha\,\mathcal{L}_{\text{CE}} + (1-\alpha)\,\mathcal{L}_{\text{HLCC}},
  \label{eq:combined-loss}
\end{equation}
where $\alpha$ decays from 1 to a target value over a warmup period.
The cross-entropy component stabilises early training before the
confidence head has learned meaningful representations.

\subsection{Uncertainty-Aware LayerNorm}

To capture the doubt signals, we extend standard LayerNorm to record
pre-normalisation statistics:

\begin{equation}
  \text{UncertaintyAwareLayerNorm}(\mathbf{x}) = \text{LayerNorm}(\mathbf{x}),
  \quad \text{storing} \quad
  \begin{bmatrix}
    \|\mathbf{x}\|_{2} \\
    \sqrt{\text{Var}(\mathbf{x}) + \epsilon} \\
    \text{Var}(\mathbf{x})
  \end{bmatrix}
\end{equation}

These three signals per layer (pre-norm magnitude, correction scale,
and variance) form the input to the norm-shift path.  For inference
with a pre-trained frozen model, only the norm-shift signals from
\texttt{output\_hidden\_states=True} are needed; no modification to the
base model is required.


% ============================================================
\section{Benchmarking Methodology}
\label{sec:methodology}
% ============================================================

\subsection{Models}

We evaluate 9 instruction-tuned models spanning 1.1B to 14.2B parameters.
Models are loaded at the highest precision that fits within the 32\,GB
VRAM budget:

\begin{table}[H]
\centering
\caption{Models evaluated with precision used for each scale.}
\label{tab:models}
\begin{tabular}{@{}llrl@{}}
\toprule
Key & Model & Parameters & Precision \\
\midrule
TinyLlama    & TinyLlama-1.1B-Chat-v1.0           & 1.1B  & FP16 \\
Qwen2-1.5B   & Qwen2-1.5B-Instruct                & 1.5B  & FP16 \\
Phi-2        & microsoft/phi-2                     & 2.7B  & FP16 \\
Phi-3-mini   & Phi-3-mini-4k-instruct              & 3.8B  & FP16 \\
Mistral-7B   & Mistral-7B-Instruct-v0.3            & 7.2B  & FP16 \\
Qwen2.5-7B   & Qwen2.5-7B-Instruct                & 7.6B  & FP16 \\
Llama-3.1-8B & Llama-3.1-8B-Instruct              & 8.0B  & FP16 \\
Qwen2.5-14B  & Qwen2.5-14B-Instruct               & 14.2B & 8-bit \\
DeepSeek-R1-14B & DeepSeek-R1-Distill-Qwen-14B     & 14.2B & 8-bit \\
\bottomrule
\end{tabular}
\end{table}

\subsection{Datasets}

We use 9 multiple-choice QA datasets covering diverse reasoning domains:

\begin{table}[H]
\centering
\caption{Evaluation datasets.  All converted to uniform MCQ format.}
\label{tab:datasets}
\begin{tabular}{@{}llcl@{}}
\toprule
Dataset & Source & Options & Domain \\
\midrule
TruthfulQA      & \cite{lin2022truthfulqa}   & 4--5 & Calibration / misconceptions \\
ARC-Easy        & \cite{clark2018arc}        & 4    & Science (elementary) \\
ARC-Challenge   & \cite{clark2018arc}        & 4    & Science (hard) \\
MMLU            & \cite{hendrycks2021mmlu}   & 4    & 57 subjects \\
HellaSwag       & \cite{zellers2019hellaswag}& 4    & Commonsense completion \\
MMLU-Pro        & \cite{wang2024mmlu}        & 10   & Extended MMLU (A--J) \\
GSM8K           & \cite{cobbe2021gsm8k}      & 4    & Math reasoning (converted) \\
BoolQ           & \cite{clark2019boolq}      & 2    & Yes/no QA \\
TriviaQA        & \cite{joshi2017triviaqa}   & 4    & Open-domain QA (converted) \\
\bottomrule
\end{tabular}
\end{table}

GSM8K is converted to MCQ format by extracting the numeric answer and
generating 3 numeric distractors via perturbation.  TriviaQA is
converted by sampling 3 distractor answers from other questions in the
dataset.  BoolQ maps directly to 2-choice MCQ.

\subsection{Confidence Methods}

We evaluate 14 confidence estimation methods (13 baselines plus our
NormShift head) organised in three tiers by computational cost, plus
two combination methods.

\paragraph{Tier 1: Single Forward Pass.}
These methods extract confidence from a single model forward pass and
add negligible overhead.

\begin{description}[nosep,leftmargin=1.5em]
  \item[Entropy] Shannon entropy of the softmax distribution over
        answer-choice tokens, normalised to $[0,1]$.
  \item[Top-$k$] Ratio of top-1 to top-2 logit probabilities.
  \item[MSP] Maximum softmax probability over choice
        tokens~\cite{hendrycks2017baseline}.
  \item[G-NLL] Greedy negative log-likelihood of the selected token.
  \item[NormShift (\hlcc{})] Trained dual-path confidence head using
        norm-shift signals and \hlcc{} loss (ours).
\end{description}

\paragraph{Tier 2: Post-Hoc Recalibration.}
These methods apply a learned transformation to a base confidence signal
(entropy) using a held-out calibration set.

\begin{description}[nosep,leftmargin=1.5em]
  \item[Temperature Scaling] Single learned temperature parameter
        $T$~\cite{guo2017calibration}.
  \item[Isotonic Regression] Non-parametric monotonic recalibration
        using the pool-adjacent-violators
        algorithm~\cite{zadrozny2002transforming}.
  \item[Platt Scaling] Two-parameter logistic regression
        $\sigma(a \cdot \text{logit}(c) + b)$~\cite{platt1999probabilistic}.
\end{description}

\paragraph{Tier 3: Multi-Pass Generation.}
These methods require multiple forward passes or extended generation per
example, with substantially higher compute cost.

\begin{description}[nosep,leftmargin=1.5em]
  \item[Verbalized] Prompt the model to self-report confidence as a
        number.
  \item[Self-Consistency] Sample $n=10$ answers and take the majority
        vote frequency as confidence~\cite{wang2023selfconsistency}.
  \item[Semantic Entropy] Cluster $n$ sampled answers by semantic
        equivalence and compute Shannon entropy over cluster
        probabilities~\cite{kuhn2023semantic}.
  \item[CoT Confidence] Two-stage prompting: generate chain-of-thought
        reasoning, then extract answer and confidence from a structured
        JSON response.
\end{description}

Additionally, we evaluate two combination methods:
\begin{description}[nosep,leftmargin=1.5em]
  \item[LVU] Linguistic Verification of Uncertainty: generate an
        explanation and detect hedging language (``might'', ``perhaps'')
        vs.\ certainty language (``definitely'', ``clearly'').
  \item[CoCoA] Confidence Combination Approach: weighted average of
        entropy and self-consistency signals.
\end{description}

\subsection{Evaluation Metrics}

We report metrics across three complementary axes.

\paragraph{Discrimination.}
Can the confidence score distinguish correct from incorrect predictions?

\begin{itemize}[nosep]
  \item \textbf{AUROC}: Area under the ROC curve, treating
        $\text{confidence} > t$ as a binary classifier for correctness.
  \item \textbf{AUPRC}: Area under the precision-recall curve, more
        informative than AUROC when the class distribution is
        imbalanced~\cite{saito2015precision}.
\end{itemize}

\paragraph{Calibration.}
Does the confidence match the empirical frequency of correctness?

\begin{itemize}[nosep]
  \item \textbf{ECE}: Expected Calibration Error with 10 equal-width
        bins~\cite{naeini2015obtaining}.
  \item \textbf{Brier Score}: Mean squared error between confidence and
        binary correctness, $\frac{1}{n}\sum(c_i - y_i)^{2}$.
  \item \textbf{Overconfidence Rate}: Fraction of incorrect predictions
        with confidence $> 0.7$.
\end{itemize}

\paragraph{Selective Prediction.}
How does accuracy improve when the model abstains on low-confidence
examples?

\begin{itemize}[nosep]
  \item \textbf{PRR}: Prediction Rejection Ratio.  Measures the area
        between the model's accuracy-vs-rejection curve and the random
        baseline, normalised by the oracle upper
        bound~\cite{malinin2017predictive}:
        \begin{equation}
          \text{PRR} = \frac{\text{AULC}_{\text{random}} -
          \text{AULC}_{\text{model}}}{\text{AULC}_{\text{random}} -
          \text{AULC}_{\text{oracle}}}.
        \end{equation}
  \item \textbf{Accuracy@Coverage}: Accuracy on the top $k$\% most
        confident predictions, reported at $k \in \{70, 80, 90\}$.
\end{itemize}

\paragraph{HLCC-Specific.}
\begin{itemize}[nosep]
  \item \textbf{HLCC Score}: Mean score under
        \Cref{eq:hlcc-scoring}, measuring confidence-weighted
        performance.
  \item \textbf{Theoretical Max HLCC}: The maximum achievable HLCC
        score for a model with accuracy $p$ is $S_{\max} = 2p$,
        attained when the model reports $c = 1$ on every correct
        prediction and $c = 0$ on every incorrect prediction (perfect
        knowledge of its own correctness).
  \item \textbf{HLCC Efficiency}: The ratio of actual to theoretical
        maximum HLCC, $\eta = S_{\text{HLCC}} / (2p)$, measuring how
        well the confidence method utilises the model's accuracy.
        A value of $\eta = 1$ indicates perfect self-knowledge.
  \item \textbf{Calibration Gap}: $|\bar{c} - \text{accuracy}|$, the
        absolute difference between mean confidence and overall accuracy.
  \item \textbf{CBM Score}: Score under the discrete CBM scheme for
        comparison with educational assessment literature.
\end{itemize}


% ============================================================
\section{Experimental Setup}
\label{sec:setup}
% ============================================================

\paragraph{Hardware.}
All experiments run on a single NVIDIA \hwGPU{} GPU
(\hwVRAM{}\,GB VRAM), \hwCPU{} CPU, and \hwRAM{}\,GB RAM.
Models of 8B parameters or fewer are loaded in FP16 precision;
14B models use 8-bit quantisation via \texttt{bitsandbytes};
models above 14B use 4-bit NF4 quantisation~\cite{dettmers2023qlora}.
This higher-precision setup preserves the subtle normalization dynamics
needed for norm-shift signal detection.

\paragraph{MCQ Evaluation Protocol.}
For each (model, dataset) pair, we evaluate up to \expCalibrationExamples{} examples.  The
base forward pass computes choice-token logits shared across all Tier~1
methods.  Tier~2 methods use a 50/50 calibration/test split.  Tier~3
methods use $n = \expSelfConSamples{}$ samples for self-consistency and semantic entropy.

\paragraph{NormShift Training.}
The confidence head is trained for \expEpochs{} epochs on MCQ data with the
combined loss from \Cref{eq:combined-loss}, using $\alpha = \expAlpha{}$ and a
\expWarmupSteps{}-step cross-entropy warmup.  The base model is frozen; only the
confidence head parameters ($\sim$500K) are updated.  We use AdamW with
learning rate $\expLR{}$ and weight decay \expWeightDecay{}.  Two variants
are trained per model:
\begin{enumerate}[nosep]
  \item \textbf{Combined}: Full dual-path architecture (norm-shift +
        hidden state).
  \item \textbf{Norm-Shift Only}: Ablation using only the norm-shift
        path, to isolate the contribution of normalisation dynamics.
\end{enumerate}

\paragraph{Multi-Day Orchestration.}
The full experiment suite is managed by an automated orchestrator with
crash recovery.  Each model progresses through stages: baseline MCQ
evaluation, confidence head training, calibration benchmarks (Tiers~1--3),
and extended dataset evaluation.  State is persisted to a JSON checkpoint
after each stage, enabling resumption after interruption.


% ============================================================
\section{Results}
\label{sec:results}
% ============================================================

\subsection{Discrimination: AUROC}

\begin{table}[H]
\centering
\caption{AUROC ($\uparrow$) and AUPRC ($\uparrow$) on TruthfulQA
         (\expCalibrationExamples{} examples).
         All values from \texttt{generate\_paper\_tables.py}.
         Best per column in \textbf{bold}.}
\label{tab:discrimination}
\small
\begin{tabular}{@{}lcccccccc@{}}
\toprule
 & \multicolumn{2}{c}{Mistral-7B} & \multicolumn{2}{c}{Qwen2.5-7B}
 & \multicolumn{2}{c}{Llama-3.1-8B} & \multicolumn{2}{c}{Qwen2.5-14B} \\
\cmidrule(lr){2-3}\cmidrule(lr){4-5}\cmidrule(lr){6-7}\cmidrule(lr){8-9}
Method & AUC & PRC & AUC & PRC & AUC & PRC & AUC & PRC \\
\midrule
Entropy          & \mistralEntropyAUROC & \mistralEntropyAUPRC
                 & \qwenEntropyAUROC    & \qwenEntropyAUPRC
                 & \llamaEntropyAUROC   & \llamaEntropyAUPRC
                 & \qwenLEntropyAUROC   & \qwenLEntropyAUPRC \\
Top-$k$          & \mistralTopkAUROC    & \mistralTopkAUPRC
                 & \qwenTopkAUROC       & \qwenTopkAUPRC
                 & \llamaTopkAUROC      & \llamaTopkAUPRC
                 & \qwenLTopkAUROC      & \qwenLTopkAUPRC \\
MSP              & \mistralMSPAUROC     & \mistralMSPAUPRC
                 & \qwenMSPAUROC        & \qwenMSPAUPRC
                 & \llamaMSPAUROC       & \llamaMSPAUPRC
                 & \qwenLMSPAUROC       & \qwenLMSPAUPRC \\
Verbalized       & \mistralVerbAUROC    & \mistralVerbAUPRC
                 & \qwenVerbAUROC       & \qwenVerbAUPRC
                 & \llamaVerbAUROC      & \llamaVerbAUPRC
                 & ---                  & --- \\
Self-Consistency & \mistralSelfConAUROC & \mistralSelfConAUPRC
                 & \qwenSelfConAUROC    & \qwenSelfConAUPRC
                 & \llamaSelfConAUROC   & \llamaSelfConAUPRC
                 & ---                  & --- \\
Temp.\ Scaling   & \mistralTempScaleAUROC & \mistralTempScaleAUPRC
                 & \qwenTempScaleAUROC  & \qwenTempScaleAUPRC
                 & \llamaTempScaleAUROC & \llamaTempScaleAUPRC
                 & ---                  & --- \\
\textbf{NormShift} & \mistralNormShiftAUROC & \mistralNormShiftAUPRC
                 & \qwenNormShiftAUROC  & \qwenNormShiftAUPRC
                 & \llamaNormShiftAUROC & \llamaNormShiftAUPRC
                 & ---                  & --- \\
\bottomrule
\end{tabular}
\end{table}

\noindent
On Mistral-7B, NormShift achieves AUROC\,=\,\mistralNormShiftAUROC{},
the highest among all methods for that model, demonstrating strong
in-domain discrimination.  On Llama-3.1-8B, NormShift reaches
AUROC\,=\,\llamaNormShiftAUROC{}, competitive with entropy
(\llamaEntropyAUROC{}).  On Qwen2.5-7B, entropy achieves the highest
AUROC (\qwenEntropyAUROC{}), while NormShift is lower
(\qwenNormShiftAUROC{}), reflecting greater architecture sensitivity.

\subsection{Calibration: ECE and Brier Score}

\begin{table}[H]
\centering
\caption{ECE ($\downarrow$) and Brier Score ($\downarrow$) on TruthfulQA
         (\expCalibrationExamples{} examples).
         All values from \texttt{generate\_paper\_tables.py}.
         Best per column in \textbf{bold}.}
\label{tab:calibration}
\small
\begin{tabular}{@{}lcccccccc@{}}
\toprule
 & \multicolumn{2}{c}{Mistral-7B} & \multicolumn{2}{c}{Qwen2.5-7B}
 & \multicolumn{2}{c}{Llama-3.1-8B} & \multicolumn{2}{c}{Qwen2.5-14B} \\
\cmidrule(lr){2-3}\cmidrule(lr){4-5}\cmidrule(lr){6-7}\cmidrule(lr){8-9}
Method & ECE & Brier & ECE & Brier & ECE & Brier & ECE & Brier \\
\midrule
Entropy          & \mistralEntropyECE   & \mistralEntropyBrier
                 & \qwenEntropyECE      & \qwenEntropyBrier
                 & \llamaEntropyECE     & \llamaEntropyBrier
                 & \qwenLEntropyECE     & \qwenLEntropyBrier \\
Top-$k$          & \mistralTopkECE      & \mistralTopkBrier
                 & \qwenTopkECE         & \qwenTopkBrier
                 & \llamaTopkECE        & \llamaTopkBrier
                 & \qwenLTopkECE        & \qwenLTopkBrier \\
MSP              & \mistralMSPECE       & \mistralMSPBrier
                 & \qwenMSPECE          & \qwenMSPBrier
                 & \llamaMSPECE         & \llamaMSPBrier
                 & \qwenLMSPECE         & \qwenLMSPBrier \\
Verbalized       & \mistralVerbECE      & \mistralVerbBrier
                 & \qwenVerbECE         & \qwenVerbBrier
                 & \llamaVerbECE        & \llamaVerbBrier
                 & ---                  & --- \\
Self-Consistency & \mistralSelfConECE   & \mistralSelfConBrier
                 & \qwenSelfConECE      & \qwenSelfConBrier
                 & \llamaSelfConECE     & \llamaSelfConBrier
                 & ---                  & --- \\
\textbf{NormShift} & \mistralNormShiftECE & \mistralNormShiftBrier
                 & \qwenNormShiftECE    & \qwenNormShiftBrier
                 & \llamaNormShiftECE   & \llamaNormShiftBrier
                 & ---                  & --- \\
\bottomrule
\end{tabular}
\end{table}

\subsection{Selective Prediction: PRR and Accuracy at Coverage}

\begin{table}[H]
\centering
\caption{PRR ($\uparrow$) and selective accuracy at coverage thresholds
         on TruthfulQA (\expCalibrationExamples{} examples).
         All values from \texttt{generate\_paper\_tables.py}.
         Best per model in \textbf{bold}.}
\label{tab:selective}
\small
\begin{tabular}{@{}lcccc@{}}
\toprule
 & \multicolumn{4}{c}{PRR ($\uparrow$)} \\
\cmidrule(lr){2-5}
Method & Mistral-7B & Qwen2.5-7B & Llama-3.1-8B & Qwen2.5-14B \\
\midrule
Entropy            & \mistralEntropyPRR   & \qwenEntropyPRR
                   & \llamaEntropyPRR     & \qwenLEntropyPRR \\
Top-$k$            & \mistralTopkPRR      & \qwenTopkPRR
                   & \llamaTopkPRR        & \qwenLTopkPRR \\
MSP                & \mistralMSPPRR       & \qwenMSPPRR
                   & \llamaMSPPRR         & \qwenLMSPPRR \\
\textbf{NormShift} & \mistralNormShiftPRR & \qwenNormShiftPRR
                   & \llamaNormShiftPRR   & --- \\
\bottomrule
\end{tabular}
\end{table}

\subsection{HLCC Score Comparison}

\begin{table}[H]
\centering
\caption{HLCC Score ($\uparrow$) on TruthfulQA (\expCalibrationExamples{} examples).
         Higher values indicate better confidence-weighted performance.
         All values from \texttt{generate\_paper\_tables.py}.
         Best per column in \textbf{bold}.}
\label{tab:hlcc-scores}
\small
\begin{tabular}{@{}lcccc@{}}
\toprule
Method & Mistral-7B & Qwen2.5-7B & Llama-3.1-8B & Qwen2.5-14B \\
\midrule
Entropy          & \mistralEntropyHLCC   & \qwenEntropyHLCC
                 & \llamaEntropyHLCC     & \qwenLEntropyHLCC \\
Top-$k$          & \mistralTopkHLCC      & \qwenTopkHLCC
                 & \llamaTopkHLCC        & \qwenLTopkHLCC \\
MSP              & \mistralMSPHLCC       & \qwenMSPHLCC
                 & \llamaMSPHLCC         & \qwenLMSPHLCC \\
Verbalized       & \mistralVerbHLCC      & \qwenVerbHLCC
                 & \llamaVerbHLCC        & --- \\
Self-Consistency & \mistralSelfConHLCC   & \qwenSelfConHLCC
                 & \llamaSelfConHLCC     & --- \\
\textbf{NormShift} & \mistralNormShiftHLCC & \qwenNormShiftHLCC
                 & \llamaNormShiftHLCC   & --- \\
\bottomrule
\end{tabular}
\end{table}

On Llama-3.1-8B, the NormShift head achieves an HLCC score of
\llamaNormShiftHLCC{}, competitive with logit-based methods.
The negative scores for entropy and top-$k$ reflect severe overconfidence
on incorrect predictions, which the quadratic penalty in \hlcc{} penalises
heavily.  Self-consistency, which uses \expSelfConSamples{}
independent samples, achieves positive scores on multiple models, but at
$10\times$ inference cost.

\subsection{Cross-Dataset Generalisation}

\begin{table}[H]
\centering
\caption{AUROC ($\uparrow$) across datasets.
         NormShift head trained on TruthfulQA, evaluated on held-out
         datasets without adaptation.  Qwen2.5-7B.
         Best per row in \textbf{bold}.}
\label{tab:cross-dataset}
\small
\begin{tabular}{@{}lcccccc@{}}
\toprule
Dataset & Acc & Entropy & MSP & Top-$k$ & NormShift & $\Delta$AUROC \\
\midrule
TruthfulQA (in-domain) & .497 & .602 & .627 & .634 & \textbf{.752} & $+$.118 \\
ARC-Challenge          & .876 & \textbf{.830} & .831 & .828 & .569 & $-$.261 \\
MMLU                   & .620 & \textbf{.792} & .778 & .761 & .468 & $-$.324 \\
HellaSwag              & .597 & .718 & \textbf{.723} & .721 & .402 & $-$.321 \\
\bottomrule
\multicolumn{7}{l}{\scriptsize $\Delta$AUROC = NormShift AUROC $-$ best logit-based AUROC.}
\end{tabular}
\end{table}

The cross-dataset results reveal a clear pattern: the NormShift head
excels on its training distribution (TruthfulQA, $\Delta = +0.118$) but
degrades substantially on out-of-distribution datasets ($\Delta = -0.26$
to $-0.32$).  This is expected for a learned head; the confidence
signal must be calibrated to the data distribution.  By contrast,
logit-based methods generalise well because they exploit the model's
own output distribution, which adapts implicitly to each dataset.

\paragraph{50/50 Split Evaluation.}
To obtain a cleaner separation of in-domain and OOD performance, we
train on 50\% of three datasets (ARC-Easy, TruthfulQA, ARC-Challenge,
150 examples each) and test on the held-out 50\% plus five unseen
datasets (\Cref{tab:cross-dataset-mistral-7b,tab:cross-dataset-qwen25-7b}).

Mistral-7B shows the strongest pattern: NormShift achieves AUROC 0.978
on held-out TruthfulQA ($+$0.215 over MSP), 0.930 on ARC-Easy
($+$0.033), and 0.919 on ARC-Challenge ($+$0.126).  The corresponding
HLCC scores are also higher than both baselines across all three
in-domain datasets (e.g., TruthfulQA: 1.125 vs.\ MSP 0.186).  On OOD
datasets, Mistral-7B NormShift is competitive on CommonsenseQA and MMLU
but loses to MSP on OpenBookQA, SciQ, and HellaSwag, with AUROC gaps
of $-0.05$ to $-0.15$.

Qwen2.5-7B tells a different story: NormShift wins on TruthfulQA
(AUROC 0.706 vs.\ MSP 0.620) but loses to MSP on the other two
in-domain datasets (ARC-Easy: 0.716 vs.\ 0.887; ARC-Challenge: 0.636
vs.\ 0.892).  On OOD data the gap widens further, with NormShift
lagging MSP by 0.15--0.33 AUROC across all five datasets.  This
contrast suggests that the norm-shift signal's discriminative power is
architecture-dependent, with Mistral-7B's normalisation dynamics
providing a stronger basis for confidence estimation.

\paragraph{Confidence-Gated Retry.}
When NormShift confidence falls below 0.4, we re-evaluate with a
trick-question system prompt.  The results are strikingly
model-dependent.  On Mistral-7B, the retry mechanism achieves a 32\%
rescue rate on TruthfulQA (25 of 78 retried questions corrected, zero
damage), producing a net accuracy gain from 0.453 to 0.620.  On
ARC-Challenge, 3 rescues against 1 damage yields a small net positive.
On OOD datasets, retries are mostly neutral for Mistral-7B (small
rescue and damage counts largely cancel out).

On Qwen2.5-7B, the retry mechanism is harmful: damage rates range from
12.5\% (TruthfulQA) to 34.7\% (OpenBookQA), with rescue rates below
4\% on most datasets.  The high retry rate (43--59\% of questions
retried on OOD data) suggests that the confidence threshold is too
aggressive for this model.  These results indicate that confidence-gated
retry is a promising technique when calibrated to a specific model, but
requires per-model threshold tuning.

\subsection{NormShift Ablation}
\label{sec:ablation}

To isolate the contribution of the norm-shift signal, we train both the
\emph{combined} (norm-shift + hidden state) and \emph{norm-shift-only}
variants for each model.  \Cref{tab:ablation} compares ECE and HLCC
scores across all five primary models.

\begin{table}[H]
\centering
\caption{Ablation: Combined vs.\ Norm-Shift-Only confidence heads.
         Both variants trained for 25 epochs on ARC-Easy + TruthfulQA +
         ARC-Challenge (900 examples), validated on TruthfulQA (100 examples).
         Results use FP16 (7--8B models) or 8-bit (14B models).
         Max HLCC\,=\,$2p$ is the theoretical maximum for a model with
         perfect self-knowledge; $\eta$\,=\,HLCC\,/\,Max HLCC is the
         efficiency.}
\label{tab:ablation}
\small
\begin{tabular}{@{}lccccccccc@{}}
\toprule
 & & & \multicolumn{3}{c}{Combined} & \multicolumn{3}{c}{Norm-Shift Only} \\
\cmidrule(lr){4-6} \cmidrule(lr){7-9}
Model & Acc & Max HLCC & ECE $\downarrow$ & HLCC $\uparrow$ & $\eta$ & ECE $\downarrow$ & HLCC $\uparrow$ & $\eta$ \\
\midrule
Mistral-7B       & 0.53 & 1.060 & \textbf{\mistralCombinedECE} & \textbf{\mistralCombinedHLCC} & 98.0\% & \mistralNSOECE & \mistralNSOHLCC & 49.0\% \\
Qwen2.5-7B       & 0.32 & 0.640 & \textbf{\qwenCombinedECE} & \textbf{\qwenCombinedHLCC} & 84.4\% & \qwenNSOECE & \qwenNSOHLCC & 51.9\% \\
Llama-3.1-8B     & 0.92 & 1.840 & \textbf{\llamaCombinedECE} & \textbf{\llamaCombinedHLCC} & 95.7\% & \llamaNSOECE & \llamaNSOHLCC & 68.8\% \\
Qwen2.5-14B      & 0.89 & 1.780 & \textbf{\qwenLCombinedECE} & \textbf{\qwenLCombinedHLCC} & 93.1\% & \qwenLNSOECE & \qwenLNSOHLCC & 80.3\% \\
DeepSeek-R1-14B  & 0.48 & 0.960 & \textbf{\deepseekCombinedECE} & \textbf{\deepseekCombinedHLCC} & 81.4\% & \deepseekNSOECE & \deepseekNSOHLCC & 4.0\% \\
\bottomrule
\end{tabular}
\end{table}

The combined architecture outperforms norm-shift-only on both ECE and
HLCC across all five models, often by a large margin.  The HLCC
efficiency metric $\eta$ contextualises raw scores by the model's
accuracy: Llama-3.1-8B has the highest raw HLCC (\llamaCombinedHLCC{})
but this is largely driven by its high accuracy (92\%), giving
$\eta = 95.7\%$.  Mistral-7B, despite its lower accuracy (53\%),
achieves the highest efficiency ($\eta = 98.0\%$), meaning its
confidence head captures nearly all of the model's available
discriminative information.

The gap between combined and norm-shift-only is most
dramatic for DeepSeek-R1-14B, where the norm-shift-only variant
essentially fails to train ($\eta = 4.0\%$ vs.\ $81.4\%$ for combined),
and for Llama-3.1-8B ($\eta = 68.8\%$ vs.\ $95.7\%$).  These results
confirm that the hidden-state path provides information that the
norm-shift signal alone cannot reliably capture.
Qwen2.5-14B is the model where the two variants are closest in HLCC
(\qwenLCombinedHLCC{} vs.\ \qwenLNSOHLCC{}, $\eta = 93.1\%$ vs.\
$80.3\%$), consistent with its strong layer-level separation signals
(\Cref{tab:layer-signals}).  Mistral-7B has the smallest ECE gap
(\mistralCombinedECE{} vs.\ \mistralNSOECE{}), but the largest
efficiency gap ($98.0\%$ vs.\ $49.0\%$), indicating that while both
variants are well-calibrated, the combined variant is far better at
assigning high confidence to correct answers.

\subsection{Per-Layer Signal Analysis}
\label{sec:layer-analysis}

To understand \emph{where} in the transformer the confidence-relevant
signal originates, we extract the standard deviation of hidden
activations at each layer for correct vs.\ incorrect predictions and
compute Cohen's $d$ as a measure of separation strength.
\Cref{tab:layer-signals} reports the top discriminative layers for each
model.

\begin{table}[H]
\centering
\caption{Per-layer norm-shift separation (Cohen's $d$) between correct
         and incorrect predictions.  90 examples per model across
         TruthfulQA, ARC-Easy, ARC-Challenge.  Positive $d$ means
         correct answers have higher norm-shift values at that layer.}
\label{tab:layer-signals}
\small
\begin{tabular}{@{}lrrcl@{}}
\toprule
Model & Layers & Acc & Best Layer & Top-3 (layer: $d$) \\
\midrule
Mistral-7B      & \mistralNLayers & \mistralLayerAcc & L\mistralBestLayer & L\mistralBestLayer: $+$\mistralBestLayerD, L\mistralSecondLayer: $+$\mistralSecondLayerD, L\mistralThirdLayer: $+$\mistralThirdLayerD \\
Qwen2.5-7B      & \qwenNLayers & \qwenLayerAcc & L\qwenBestLayer & L\qwenBestLayer: $+$\qwenBestLayerD, L\qwenSecondLayer: $+$\qwenSecondLayerD, L\qwenThirdLayer: $-$\qwenThirdLayerD \\
Llama-3.1-8B    & \llamaNLayers & \llamaLayerAcc & L\llamaBestLayer  & L\llamaBestLayer: $+$\llamaBestLayerD, L\llamaSecondLayer: $-$\llamaSecondLayerD, L\llamaThirdLayer: $-$\llamaThirdLayerD \\
Qwen2.5-14B     & \qwenLNLayers & \qwenLayerAcc & L\qwenLBestLayer & L\qwenLBestLayer: $-$\qwenLBestLayerD, L\qwenLSecondLayer: $+$\qwenLSecondLayerD, L\qwenLThirdLayer: $-$\qwenLThirdLayerD \\
DeepSeek-R1-14B & \deepseekNLayers & \deepseekLayerAcc & L\deepseekBestLayer & L\deepseekBestLayer: $-$\deepseekBestLayerD, L\deepseekSecondLayer: $-$\deepseekSecondLayerD, L\deepseekThirdLayer: $+$\deepseekThirdLayerD \\
\bottomrule
\end{tabular}
\end{table}

Several patterns emerge from the FP16 extraction results across five
models.  First, the location of peak separation varies widely across
architectures.  Qwen2.5-14B concentrates its signal in the middle
layers (L\qwenLBestLayer{}, $|d| = \qwenLBestLayerD{}$), while Mistral-7B peaks in the upper layers
(L\mistralBestLayer{}, $d = +\mistralBestLayerD{}$).  Most surprisingly, Llama-3.1-8B shows its
strongest separation at Layer~\llamaBestLayer{} ($d = +\llamaBestLayerD{}$), suggesting that for this
architecture the residual stream carries confidence-relevant structure
from very early in the forward pass.  Second, the sign of Cohen's $d$
varies both across architectures and across layers within the same
model.  Qwen2.5-14B shows negative $d$ at L30 (correct answers have
\emph{lower} norm-shift) but positive $d$ at L47, while Llama-3.1-8B
has positive $d$ at L2 and negative $d$ at L3 and L14.  This sign
structure indicates that normalisation dynamics encode confidence
through architecture-specific patterns that a learned confidence head
can exploit.  Third, DeepSeek-R1-14B shows the weakest peak separation
(max $|d| = \deepseekBestLayerD{}$), consistent with its higher ECE in the ablation
results (\Cref{tab:ablation}).  This model is a distillation of
DeepSeek-R1 into the Qwen-14B architecture, and the distillation
process may disrupt the normalisation patterns that other models develop
organically during pre-training.

\subsection{Theoretical Analysis}

\Cref{fig:theory} shows the theoretical optimal confidence curves
derived in \Cref{sec:theory}.  The \hlcc{} curve lies strictly below
the Brier identity line for $p < 0.75$, confirming the conservative
bias.  The CBM discrete scheme approximates \hlcc{} with staircase
steps.

\begin{figure}[H]
\centering
\begin{subfigure}[b]{0.48\textwidth}
  \includegraphics[width=\textwidth]{../data/results/theory/optimal_confidence_comparison.png}
  \caption{Optimal confidence vs.\ true accuracy.}
  \label{fig:optimal-conf}
\end{subfigure}
\hfill
\begin{subfigure}[b]{0.48\textwidth}
  \includegraphics[width=\textwidth]{../data/results/theory/reward_penalty_asymmetry.png}
  \caption{Reward/penalty asymmetry.}
  \label{fig:asymmetry}
\end{subfigure}
\caption{Theoretical properties of the \hlcc{} scoring function.
         (a)~The optimal confidence under \hlcc{} is conservative
         ($\cstar < p$) below $p = 0.75$.
         (b)~The linear reward $R(c) = 1+c$ vs.\ quadratic penalty
         $P(c) = -2c^{2}$ creates asymmetric marginal costs.}
\label{fig:theory}
\end{figure}


% ============================================================
\section{Discussion}
\label{sec:discussion}
% ============================================================

\paragraph{Conservative Bias as a Feature.}
The \hlcc{} optimal confidence formula $\cstar = p/[4(1-p)]$
systematically underestimates confidence for models with accuracy below
75\%.  While this departs from the strict propriety of Brier and
log-loss, it may be advantageous in deployment scenarios where
overconfident errors carry disproportionate cost, for example in medical
diagnosis, legal reasoning, or autonomous decision-making.

\paragraph{Model-Dependent NormShift Effectiveness.}
The NormShift head exhibits striking performance variation across model
architectures and scales.  Among five primary models trained with
identical hyperparameters at FP16 or 8-bit precision, the
\emph{combined} variant achieves its best calibration on Mistral-7B
(ECE\,=\,\mistralCombinedECE{}), followed by Llama-3.1-8B (ECE\,=\,\llamaCombinedECE{}) and
Qwen2.5-7B (ECE\,=\,\qwenCombinedECE{}).  Llama-3.1-8B attains the highest HLCC
score (\llamaCombinedHLCC{}) driven by its high validation accuracy, closely followed
by Qwen2.5-14B (HLCC\,=\,\qwenLCombinedHLCC{}).  The 14B models show
that scaling does not uniformly improve calibration; DeepSeek-R1-14B
achieves ECE\,=\,\deepseekCombinedECE{}, considerably worse than Qwen2.5-14B
(ECE\,=\,\qwenLCombinedECE{}) despite identical parameter counts.  This variation
suggests that the utility of norm-shift signals depends on the model's
internal normalisation dynamics and residual stream geometry, not merely
on model scale.  DeepSeek-R1-14B, a distillation product, shows the
weakest norm-shift separation (peak $|d| = \deepseekBestLayerD{}$) and the poorest
norm-shift-only training outcome (HLCC\,=\,\deepseekNSOHLCC{}).

\paragraph{NormShift Signal Interpretation.}
The norm-shift signal $\delta^{(\ell)}$ captures how much each
transformer layer's normalisation step altered the representation.
Layers that apply large corrections may indicate inputs that fall outside
the model's learned distribution, providing a complementary uncertainty
signal to logit-based methods.  The ablation (\Cref{tab:ablation})
quantifies the contribution of this signal relative to the standard
hidden-state approach.

\paragraph{Computational Cost Tradeoffs.}
Tier~1 methods (entropy, MSP, NormShift) add negligible cost to
inference.  The NormShift head requires a single additional forward pass
through $\sim$500K parameters (less than 0.01\% of the base model).
In contrast, Tier~3 methods (self-consistency, semantic entropy) multiply
inference cost by the sample count $n$.  Across all five primary models,
the combined NormShift head achieves competitive or superior calibration
to multi-pass methods at a fraction of the cost, making it practical for
real-time confidence estimation in deployment scenarios, including
interactive token-level confidence visualisation during generation.

\paragraph{Limitations.}
\begin{enumerate}[nosep]
  \item Models above 8B parameters are evaluated under 8-bit quantisation;
        full FP16 results may differ slightly.
  \item The MCQ evaluation format may not fully capture open-ended
        generation confidence.
  \item The NormShift head is trained per-model and does not transfer
        across architectures.
  \item Performance varies substantially across model architectures,
        suggesting architecture-dependent factors in norm-shift signal
        quality.
  \item The HLCC scoring function's conservative bias may undervalue
        confidence in well-calibrated high-accuracy models.
\end{enumerate}


% ============================================================
\section{LM-Polygraph Integration}
\label{sec:polygraph}
% ============================================================

To enable direct comparison with the broader uncertainty estimation
literature, we implement the NormShift confidence head as a custom
estimator within the
LM-Polygraph~\cite{fadeeva2024lmpolygraph} benchmarking framework.
LM-Polygraph provides over 40 uncertainty estimation methods and
standardised evaluation protocols, making it the primary benchmark
for comparing UE methods across models and tasks.

\subsection{Integration Architecture}

The integration wraps the NormShift confidence head as an LM-Polygraph
\texttt{Estimator} subclass.  The key challenge is that NormShift
requires full hidden states from all transformer layers
(\texttt{output\_hidden\_states=True}), while most LM-Polygraph methods
operate only on output logits.

We implement a custom \texttt{StatCalculator} that extracts three
signals from each forward pass:
\begin{enumerate}[nosep]
  \item The full hidden states tuple
        $({\mathbf{h}^{(0)}, \ldots, \mathbf{h}^{(L)}})$,
  \item Per-layer norm-shift signals
        $\boldsymbol{\delta} \in \R^{L}$ via
        $\delta^{(\ell)} = 1 - \text{std}(\mathbf{h}^{(\ell)})$,
  \item The final hidden state $\mathbf{h}^{(L)}$ for the combined path.
\end{enumerate}

The estimator inverts the confidence output to match LM-Polygraph's
convention (higher values indicate greater uncertainty):
\begin{equation}
  u_{\text{NormShift}} = 1 - c_{\text{NormShift}}.
\end{equation}

Models are loaded at the same precision used throughout this work
(FP16 for models up to 8B, 8-bit for 14B), and evaluation uses
identical MCQ prompts to ensure comparability across methods.

\subsection{LM-Polygraph Benchmark Results}

\Cref{tab:polygraph} presents the initial benchmark comparison on
TruthfulQA using Qwen2.5-7B.  Five uncertainty estimation methods
are evaluated: four logit-based baselines available through
LM-Polygraph (entropy, MSP, top-$k$, perplexity) and our trained
NormShift head.

\begin{table}[H]
\centering
\caption{LM-Polygraph benchmark (pre-retrain): Qwen2.5-7B on TruthfulQA
         (817 examples).  NormShift uses ARC-Easy-trained checkpoint.
         Best per column in \textbf{bold}.}
\label{tab:polygraph}
\small
\begin{tabular}{@{}lcccccccc@{}}
\toprule
Method & Acc & Conf & Gap$\downarrow$ & ECE$\downarrow$ & Brier$\downarrow$
       & AUROC$\uparrow$ & AUPRC$\uparrow$ & PRR$\uparrow$ \\
\midrule
Entropy      & 0.497 & 0.578 & \textbf{0.081} & \textbf{0.145} & \textbf{0.278}
             & 0.602 & 0.539 & 0.097 \\
MSP          & 0.497 & 0.759 & 0.262 & 0.262 & 0.309
             & 0.627 & 0.559 & 0.157 \\
Top-$k$      & 0.497 & 0.819 & 0.322 & 0.322 & 0.339
             & \textbf{0.634} & \textbf{0.569} & \textbf{0.177} \\
Perplexity   & 0.497 & 0.759 & 0.262 & 0.262 & 0.309
             & 0.627 & 0.559 & 0.157 \\
NormShift$^{\dagger}$ & 0.497 & 0.821 & 0.324 & 0.330 & 0.368
             & 0.416 & 0.424 & $-$0.258 \\
\bottomrule
\multicolumn{9}{l}{\scriptsize $\dagger$ Out-of-distribution: checkpoint trained on ARC-Easy.
         See \Cref{tab:polygraph-retrained} for in-domain results.}
\end{tabular}
\end{table}

\paragraph{Analysis.}
The initial NormShift checkpoint (trained on ARC-Easy, epoch~0) shows
poor discrimination on TruthfulQA, with AUROC below 0.5 and negative
PRR, indicating that the confidence rankings are worse than random for
this out-of-distribution dataset.  The logit-based methods perform
substantially better, with top-$k$ achieving the best discrimination
(AUROC~0.634, PRR~0.177) and entropy achieving the best calibration
(ECE~0.145).

This result motivates \emph{in-domain retraining}: the NormShift head
must be trained on data from the target distribution to learn
meaningful confidence estimates.  \Cref{tab:polygraph-retrained} shows
results after retraining on TruthfulQA (600 training examples, 15
epochs, best checkpoint at epoch~1 by validation HLCC).

\paragraph{Retrained Results.}
After in-domain training, NormShift achieves a substantial improvement:
AUROC rises from 0.416 to \textbf{0.752} ($+$0.336), surpassing all
logit-based methods (next best: top-$k$ at 0.634).  PRR jumps from
$-$0.258 to \textbf{0.484}, indicating strong selective prediction: the
model can reliably reject uncertain predictions.  AUPRC reaches
\textbf{0.701} vs.\ 0.569 for top-$k$.

The mean confidence drops to 0.256 (from 0.821 pre-retrain), reflecting
the conservative bias of \hlcc{} training.  While this increases the
calibration gap (0.241) relative to entropy (0.081), the discrimination
and selective prediction gains are substantial.  Notably, the Brier
score (0.289) is competitive with entropy (0.278), confirming that the
overall probabilistic quality remains strong.

\begin{table}[H]
\centering
\caption{LM-Polygraph benchmark after NormShift retraining on TruthfulQA
         (600 train examples, 15 epochs).
         Qwen2.5-7B, 817 evaluation examples.
         Best per column in \textbf{bold}.}
\label{tab:polygraph-retrained}
\small
\begin{tabular}{@{}lcccccccc@{}}
\toprule
Method & Acc & Conf & Gap$\downarrow$ & ECE$\downarrow$ & Brier$\downarrow$
       & AUROC$\uparrow$ & AUPRC$\uparrow$ & PRR$\uparrow$ \\
\midrule
Entropy      & 0.497 & 0.578 & \textbf{0.081} & \textbf{0.145} & \textbf{0.278}
             & 0.602 & 0.539 & 0.097 \\
MSP          & 0.497 & 0.759 & 0.262 & 0.262 & 0.309
             & 0.627 & 0.559 & 0.157 \\
Top-$k$      & 0.497 & 0.819 & 0.322 & 0.322 & 0.339
             & 0.634 & 0.569 & 0.177 \\
\textbf{NormShift} & 0.497 & 0.256 & 0.241 & 0.254 & 0.289
             & \textbf{0.752} & \textbf{0.701} & \textbf{0.484} \\
\bottomrule
\end{tabular}
\end{table}

\paragraph{Cross-Dataset Transfer.}
\Cref{tab:polygraph-transfer} reports the full metric comparison when
the TruthfulQA-trained head is evaluated on held-out datasets.

\begin{table}[H]
\centering
\caption{NormShift cross-dataset transfer: trained on TruthfulQA,
         evaluated elsewhere.  Qwen2.5-7B.  PRR ($\uparrow$) and
         AUROC ($\uparrow$) shown for each method.}
\label{tab:polygraph-transfer}
\small
\begin{tabular}{@{}lcc|cc|cc|cc@{}}
\toprule
 & \multicolumn{2}{c|}{Entropy} & \multicolumn{2}{c|}{Top-$k$}
 & \multicolumn{2}{c|}{NormShift} & \multicolumn{2}{c}{Best logit} \\
\cmidrule(lr){2-3}\cmidrule(lr){4-5}\cmidrule(lr){6-7}\cmidrule(lr){8-9}
Dataset & AUC & PRR & AUC & PRR & AUC & PRR & AUC & PRR \\
\midrule
TruthfulQA     & .602 & .097 & .634 & .177 & \textbf{.752} & \textbf{.484}
               & .634 & .177 \\
ARC-Challenge  & .830 & .771 & .828 & .768 & .569 & .277
               & \textbf{.831} & \textbf{.772} \\
MMLU           & \textbf{.792} & \textbf{.679} & .761 & .641 & .468 & $-$.090
               & .792 & .679 \\
HellaSwag      & .718 & .552 & .721 & .547 & .402 & $-$.182
               & \textbf{.723} & \textbf{.553} \\
\bottomrule
\end{tabular}
\end{table}


% ============================================================
\section{Related Work}
\label{sec:related}
% ============================================================

\paragraph{Confidence Calibration in Neural Networks.}
\citet{guo2017calibration} demonstrated that modern neural networks are
poorly calibrated, with temperature scaling as a simple post-hoc fix.
\citet{naeini2015obtaining} introduced ECE as the standard calibration
metric.  \citet{nixon2019measuring} showed that ECE is sensitive to
binning choices and proposed adaptive calibration error.

\paragraph{Uncertainty Estimation in LLMs.}
\citet{kadavath2022language} studied self-evaluation in large language
models, finding that models can estimate their own accuracy.
\citet{kuhn2023semantic} proposed semantic entropy as a
generation-aware uncertainty measure.  \citet{xiong2024llms} provided a
comprehensive survey of LLM confidence estimation, categorising methods
into logit-based, verbalised, and consistency-based
approaches.  \citet{lin2024generating} and \citet{tian2023just} explored
verbalized confidence elicitation and found that prompting strategies
significantly affect self-reported confidence quality.

\paragraph{Selective Prediction.}
\citet{geifman2017selective} formalised selective prediction as
accuracy-coverage tradeoffs.  The prediction rejection ratio (PRR) was
introduced by \citet{malinin2017predictive} and adopted by the
LM-Polygraph framework~\cite{fadeeva2024lmpolygraph} as a standard
metric for uncertainty benchmarking.

\paragraph{Confidence-Based Marking.}
\citet{gardnermedwin1995} introduced Certainty-Based Marking (CBM) for
educational assessment, where students declare confidence on a
three-level scale: $C{=}1$ (low), $C{=}2$ (mid), $C{=}3$ (high).
The scoring table awards $+1/0$, $+2/{-}2$, $+3/{-}6$ for
correct/incorrect at each level, creating a discrete penalty structure
whose thresholds correspond to proper scoring boundaries at $p = 2/3$
and $p = 4/5$~\cite{gardnermedwin2006}.  This scheme is a discretised
form of a logarithmic proper scoring rule~\cite{good1952,dawid1986},
where the penalty for confident errors grows faster than the reward for
confident correctness.  \citet{gardnermedwin2003} demonstrated that CBM
improves learning engagement in computer-assisted assessment, and
\citet{gardnermedwin2019} provided evidence that certainty-based marking
stimulates deeper thinking in objective tests.
\hlcc{} can be viewed as a continuous generalisation of CBM: where CBM
discretises confidence into three levels, \hlcc{} allows any $c \in [0,1]$
while preserving the asymmetric penalty structure.  The CBM boundary at
$p = 4/5$ (above which maximum confidence is optimal) corresponds to
the \hlcc{} saturation point $\cstar = 1$ at $p = 0.8$
(\Cref{cor:crossover}).  \citet{budescu2012optimal} studied optimal
scoring rules for confidence-weighted testing, analysing the tradeoff
between strict propriety and behavioural desiderata such as participation
incentives.


% ============================================================
\section{Conclusion}
\label{sec:conclusion}
% ============================================================

We have presented the \hlcc{} scoring function and the NormShift
confidence head as a practical method for training calibrated confidence
estimates in language models.  Across five primary models spanning 7B to
14B parameters, our key findings are:

\begin{enumerate}[nosep]
  \item \textbf{In-domain dominance on Mistral-7B.}  When trained on the
        target distribution, NormShift achieves near-perfect
        discrimination on Mistral-7B (AUROC\,=\,\mistralNormShiftAUROC{},
        PRR\,=\,\mistralNormShiftPRR{} on TruthfulQA) while requiring
        only a single forward pass through $\sim$500K additional
        parameters.  The 50/50 split cross-dataset evaluation
        (\Cref{tab:cross-dataset-mistral-7b}) confirms this: Mistral-7B
        NormShift achieves AUROC\,=\,0.978 on held-out TruthfulQA,
        0.930 on ARC-Easy, and 0.919 on ARC-Challenge, surpassing MSP
        on all three in-domain datasets.  This is the strongest
        in-domain result among the models we tested.
  \item \textbf{Effective confidence-gated retry on Mistral-7B.}  When
        NormShift confidence falls below 0.4, re-evaluating with a
        trick-question prompt rescues 25 of 78 low-confidence questions
        on TruthfulQA (32\% rescue rate) with zero damage.  This
        demonstrates that NormShift can reliably identify questions worth
        retrying.  However, this mechanism is model-dependent:
        Qwen2.5-7B shows high damage rates (up to 35\% on OpenBookQA),
        indicating that the retry strategy requires per-model tuning.
  \item \textbf{Strong calibration at scale.}  The combined NormShift
        head achieves ECE\,=\,\mistralCombinedECE{} on Mistral-7B and ECE\,=\,\llamaCombinedECE{}
        on Llama-3.1-8B at FP16 precision.  Llama-3.1-8B achieves the
        highest HLCC score (\llamaCombinedHLCC{}) driven by its high
        accuracy on the training distribution.
  \item \textbf{Architecture sensitivity.}  Performance varies
        substantially across models and the strong Mistral-7B results
        do not generalise to all architectures.  On Qwen2.5-7B and
        Llama-3.1-8B, NormShift is outperformed by multi-pass baselines
        (self-consistency, semantic entropy) and even by MSP in
        cross-dataset evaluation.  DeepSeek-R1-14B shows weaker
        calibration (ECE\,=\,\deepseekCombinedECE{}) despite identical
        parameter count to Qwen2.5-14B (ECE\,=\,\qwenLCombinedECE{}).
        The norm-shift-only ablation highlights this further: DeepSeek
        essentially fails to train
        (HLCC\,=\,\deepseekNSOHLCC{}) while Qwen2.5-14B reaches \qwenLNSOHLCC{}.
  \item \textbf{Limited transfer.}  Cross-dataset evaluation reveals
        that NormShift degrades on out-of-distribution data.  On
        Mistral-7B, $\Delta$AUROC ranges from $-0.03$ (SciQ) to
        $-0.15$ (HellaSwag) on OOD datasets, with MMLU being the one
        exception where NormShift narrowly leads ($+0.01$).  On
        Qwen2.5-7B the degradation is more severe ($\Delta$AUROC
        $= -0.15$ to $-0.33$), and NormShift loses to MSP on all OOD
        datasets.  Logit-based methods generalise better because they
        exploit the model's own output distribution.
  \item \textbf{Conservative bias as design choice.}  The \hlcc{}
        scoring function induces a conservative confidence mapping
        ($\cstar < p$ for $p < 0.75$), which heavily penalises
        overconfident errors.  On models where entropy and top-$k$
        produce negative HLCC scores due to overconfidence (e.g.,
        Llama-3.1-8B: entropy $-0.211$, top-$k$ $-0.314$), NormShift
        still achieves a positive score (0.324), though MSP (0.144)
        also remains positive on this model.
\end{enumerate}

The theoretical analysis confirms that \hlcc{} departs from strict
propriety in favour of asymmetric risk: rewarding correct confidence
linearly while penalising incorrect confidence quadratically.

The strongest practical result is Mistral-7B, where NormShift serves
as both a high-quality confidence estimator and an effective retry gate
for the target distribution.  For other architectures and for
out-of-distribution inputs, logit-based methods (particularly MSP and
entropy) remain more reliable, and future work should investigate
domain-agnostic training strategies for the confidence head.
We provide an interactive confidence visualisation tool that colours
generated tokens by their confidence level, enabling real-time
inspection of model certainty during inference.


% ============================================================
% APPENDIX
% ============================================================
\appendix

\section{Full Calibration Scores by Method}
\label{app:calibration-scores}

This appendix reports the complete per-method calibration scores
evaluated on TruthfulQA (300 examples).  These results complement
the summary comparisons in \Cref{sec:baselines} by providing the
full metric suite for each method and model.

\subsection{Mistral-7B (FP16)}
\label{app:calibration-mistral-7b}

\begin{table}[H]
\centering
\caption{Calibration scores for Mistral-7B on TruthfulQA (300 examples, FP16).}
\label{tab:full-calibration-mistral-7b}
\small
\begin{tabular}{@{}lccccccc@{}}
\toprule
Method & Acc & ECE\,$\downarrow$ & Brier\,$\downarrow$ & HLCC\,$\uparrow$
       & AUROC\,$\uparrow$ & AUPRC\,$\uparrow$ & PRR\,$\uparrow$ \\
\midrule
\multicolumn{8}{l}{\emph{Phase 1: Logit-based}} \\
Entropy                  & .467 & .389 & .393 & .105
                 & .602 & .531 & .159 \\
Top-$k$                  & .467 & .428 & .422 & .060
                 & .602 & .545 & .194 \\
MSP                      & .467 & .353 & .352 & .214
                 & .711 & .706 & .513 \\
GNLL                     & .467 & .353 & .352 & .214
                 & .711 & .706 & .513 \\
\midrule
\multicolumn{8}{l}{\emph{Phase 2: Generation-based}} \\
Verbalized               & .467 & .502 & .516 & $-$0.075
                 & .595 & .557 & .222 \\
Self-consistency         & .290 & .422 & .410 & $-$0.254
                 & .556 & .338 & .114 \\
CoT confidence           & .513 & .463 & .449 & .124
                 & .581 & .601 & .214 \\
Semantic entropy         & .300 & .328 & .354 & $-$0.076
                 & .518 & .323 & .054 \\
LVU                      & .467 & .122 & .287 & .412
                 & .464 & .453 & $-$0.044 \\
CoCoA                    & .277 & .506 & .456 & $-$0.385
                 & .585 & .323 & .115 \\
\midrule
\multicolumn{8}{l}{\emph{Phase 2: Post-hoc calibration$^\dagger$}} \\
Temp.\ scaling           & .467 & .515 & .510 & $-$0.090
                 & .651 & .641 & .393 \\
Isotonic$^\dagger$       & .467 & \textbf{.000} & .230 & .485
                 & .654 & .627 & .368 \\
Platt$^\dagger$          & .467 & .050 & .243 & .462
                 & .602 & .531 & .159 \\
\midrule
\multicolumn{8}{l}{\emph{Learned: NormShift (combined)$^\ddagger$}} \\
\textbf{NormShift}$^\ddagger$ & .467 & .030 & \textbf{.025} & \textbf{.884}
                 & \textbf{.999} & \textbf{.999} & \textbf{.999} \\
\bottomrule
\multicolumn{8}{l}{\scriptsize $^\dagger$Post-hoc methods fit and evaluate on the same data.} \\
\multicolumn{8}{l}{\scriptsize $^\ddagger$NormShift includes TruthfulQA in training data.}
\end{tabular}
\end{table}

\subsection{Qwen2.5-7B (FP16)}
\label{app:calibration-qwen2.5-7b}

\begin{table}[H]
\centering
\caption{Calibration scores for Qwen2.5-7B on TruthfulQA (300 examples, FP16).}
\label{tab:full-calibration-qwen2.5-7b}
\small
\begin{tabular}{@{}lccccccc@{}}
\toprule
Method & Acc & ECE\,$\downarrow$ & Brier\,$\downarrow$ & HLCC\,$\uparrow$
       & AUROC\,$\uparrow$ & AUPRC\,$\uparrow$ & PRR\,$\uparrow$ \\
\midrule
\multicolumn{8}{l}{\emph{Phase 1: Logit-based}} \\
Entropy                  & .493 & .449 & .425 & .125
                 & .899 & .895 & .834 \\
Top-$k$                  & .493 & .476 & .465 & .050
                 & .757 & .769 & .601 \\
MSP                      & .493 & .277 & .310 & .340
                 & .657 & .589 & .228 \\
GNLL                     & .493 & .277 & .310 & .340
                 & .657 & .589 & .228 \\
\midrule
\multicolumn{8}{l}{\emph{Phase 2: Generation-based}} \\
Verbalized               & .493 & .311 & .404 & .399
                 & .384 & .419 & $-$0.253 \\
Self-consistency         & .683 & .269 & .268 & \textbf{.825}
                 & .626 & .754 & .254 \\
CoT confidence           & .657 & .262 & .277 & .726
                 & .637 & .752 & .312 \\
Semantic entropy         & .687 & .261 & .277 & .824
                 & .602 & .739 & .193 \\
LVU                      & .493 & .130 & .285 & .518
                 & .554 & .510 & .037 \\
CoCoA                    & .677 & .266 & .267 & .801
                 & .806 & .904 & .726 \\
\midrule
\multicolumn{8}{l}{\emph{Phase 2: Post-hoc calibration$^\dagger$}} \\
Temp.\ scaling           & .493 & .487 & .487 & .009
                 & .652 & .586 & .235 \\
Isotonic$^\dagger$       & .493 & \textbf{.000} & \textbf{.119} & .746
                 & \textbf{.909} & \textbf{.912} & \textbf{.860} \\
Platt$^\dagger$          & .493 & .037 & .130 & .731
                 & .899 & .895 & .834 \\
\midrule
\multicolumn{8}{l}{\emph{Learned: NormShift (combined)$^\ddagger$}} \\
\textbf{NormShift}$^\ddagger$ & .493 & .305 & .339 & .460
                 & .633 & .574 & .211 \\
\bottomrule
\multicolumn{8}{l}{\scriptsize $^\dagger$Post-hoc methods fit and evaluate on the same data.} \\
\multicolumn{8}{l}{\scriptsize $^\ddagger$NormShift includes TruthfulQA in training data.}
\end{tabular}
\end{table}

\subsection{Llama-3.1-8B (FP16)}
\label{app:calibration-llama3.1-8b}

\begin{table}[H]
\centering
\caption{Calibration scores for Llama-3.1-8B on TruthfulQA (300 examples, FP16).}
\label{tab:full-calibration-llama3.1-8b}
\small
\begin{tabular}{@{}lccccccc@{}}
\toprule
Method & Acc & ECE\,$\downarrow$ & Brier\,$\downarrow$ & HLCC\,$\uparrow$
       & AUROC\,$\uparrow$ & AUPRC\,$\uparrow$ & PRR\,$\uparrow$ \\
\midrule
\multicolumn{8}{l}{\emph{Phase 1: Logit-based}} \\
Entropy                  & .393 & .524 & .491 & $-$0.211
                 & .809 & .741 & .653 \\
Top-$k$                  & .393 & .569 & .548 & $-$0.314
                 & .746 & .626 & .481 \\
MSP                      & .393 & .295 & .312 & .144
                 & .658 & .527 & .302 \\
GNLL                     & .393 & .295 & .312 & .144
                 & .658 & .527 & .302 \\
\midrule
\multicolumn{8}{l}{\emph{Phase 2: Generation-based}} \\
Verbalized               & .393 & .421 & .565 & $-$0.198
                 & .467 & .371 & $-$0.069 \\
Self-consistency         & .507 & .242 & .279 & .440
                 & .707 & .716 & .485 \\
CoT confidence           & .647 & .213 & .261 & \textbf{.744}
                 & .630 & .741 & .293 \\
Semantic entropy         & .507 & .205 & .253 & .526
                 & .691 & .744 & .527 \\
LVU                      & .393 & .120 & .251 & .314
                 & .545 & .459 & .143 \\
CoCoA                    & .523 & .301 & .309 & .393
                 & .722 & \textbf{.777} & .577 \\
\midrule
\multicolumn{8}{l}{\emph{Phase 2: Post-hoc calibration$^\dagger$}} \\
Temp.\ scaling           & .393 & .566 & .558 & $-$0.333
                 & .636 & .485 & .217 \\
Isotonic$^\dagger$       & .393 & \textbf{.000} & \textbf{.161} & .481
                 & \textbf{.828} & .751 & \textbf{.668} \\
Platt$^\dagger$          & .393 & .075 & .177 & .469
                 & .809 & .741 & .653 \\
\midrule
\multicolumn{8}{l}{\emph{Learned: NormShift (combined)$^\ddagger$}} \\
\textbf{NormShift}$^\ddagger$ & .393 & .235 & .253 & .324
                 & .794 & .713 & .613 \\
\bottomrule
\multicolumn{8}{l}{\scriptsize $^\dagger$Post-hoc methods fit and evaluate on the same data.} \\
\multicolumn{8}{l}{\scriptsize $^\ddagger$NormShift includes TruthfulQA in training data.}
\end{tabular}
\end{table}

\subsection{Qwen2.5-14B (8-bit)}
\label{app:calibration-qwen2.5-14b}

\begin{table}[H]
\centering
\caption{Calibration scores for Qwen2.5-14B on TruthfulQA (300 examples, 8-bit).}
\label{tab:full-calibration-qwen2.5-14b}
\small
\begin{tabular}{@{}lccccccc@{}}
\toprule
Method & Acc & ECE\,$\downarrow$ & Brier\,$\downarrow$ & HLCC\,$\uparrow$
       & AUROC\,$\uparrow$ & AUPRC\,$\uparrow$ & PRR\,$\uparrow$ \\
\midrule
\multicolumn{8}{l}{\emph{Phase 1: Logit-based}} \\
Entropy                  & .620 & .322 & .311 & .600
                 & .901 & .932 & .845 \\
Top-$k$                  & .620 & .352 & .342 & .549
                 & .799 & .838 & .624 \\
MSP                      & .620 & .240 & .263 & .694
                 & .719 & .796 & .508 \\
GNLL                     & .620 & .240 & .263 & .694
                 & .719 & .796 & .508 \\
\midrule
\multicolumn{8}{l}{\emph{Phase 2: Generation-based}} \\
Verbalized               & .620 & .229 & .298 & .731
                 & .604 & .707 & .264 \\
Self-consistency         & .740 & .178 & .180 & 1.113
                 & .747 & .859 & .490 \\
Semantic entropy         & .747 & .171 & .177 & \textbf{1.140}
                 & .740 & .876 & .542 \\
\midrule
\multicolumn{8}{l}{\emph{Phase 2: Post-hoc calibration$^\dagger$}} \\
Temp.\ scaling           & .620 & .366 & .361 & .516
                 & .669 & .762 & .420 \\
Isotonic$^\dagger$       & .620 & \textbf{.000} & \textbf{.112} & 1.007
                 & \textbf{.910} & \textbf{.941} & \textbf{.867} \\
Platt$^\dagger$          & .620 & .057 & .122 & .984
                 & .901 & .932 & .845 \\
\midrule
\multicolumn{8}{l}{\emph{Learned: NormShift (combined)$^\ddagger$}} \\
\textbf{NormShift}$^\ddagger$ & .620 & .183 & .201 & .840
                 & .786 & .811 & .558 \\
\bottomrule
\multicolumn{8}{l}{\scriptsize $^\dagger$Post-hoc methods fit and evaluate on the same data.} \\
\multicolumn{8}{l}{\scriptsize $^\ddagger$NormShift includes TruthfulQA in training data.}
\end{tabular}
\end{table}


\section{Cross-Dataset Evaluation (50/50 Split)}
\label{app:cross-dataset-eval}

Each training dataset (ARC-Easy, TruthfulQA, ARC-Challenge) is
split 50/50: the first half trains the NormShift head, the second
half serves as the in-domain test set.  Out-of-distribution (OOD)
datasets are never seen during training.  When NormShift confidence
falls below 0.4, a retry pass re-evaluates the question with a
trick-question system prompt.

\subsection{Mistral-7B}

\begin{table}[H]
\centering
\caption{Cross-dataset evaluation for Mistral-7B: NormShift (with trick-question retry) vs.\ logit baselines.
         Trained on 50\% of ARC-Easy + TruthfulQA + ARC-Challenge,
         tested on held-out 50\% (in-domain) and unseen datasets (OOD).}
\label{tab:cross-dataset-mistral-7b}
\small
\begin{tabular}{@{}llcccccc@{}}
\toprule
Dataset & Type & Acc & NS\,AUROC & MSP\,AUROC & $\Delta$ & NS\,HLCC & MSP\,HLCC \\
\midrule
Arc Easy             & in-dom & .900 & \textbf{.930} & .897 & $+$0.033 & 1.751 & 1.651 \\
Truthfulqa           & in-dom & .620 & \textbf{.978} & .763 & $+$0.215 & 1.125 & .186 \\
Arc Challenge        & in-dom & .760 & \textbf{.919} & .793 & $+$0.126 & 1.321 & 1.091 \\
Commonsenseqa        & OOD & .687 & .713 & \textbf{.764} & $-$0.051 & .926 & .919 \\
Openbookqa           & OOD & .747 & .644 & \textbf{.740} & $-$0.096 & 1.013 & 1.070 \\
Sciq                 & OOD & .860 & .801 & \textbf{.834} & $-$0.033 & 1.479 & 1.534 \\
Hellaswag            & OOD & .560 & .493 & \textbf{.641} & $-$0.148 & .312 & .541 \\
Mmlu                 & OOD & .413 & \textbf{.649} & .638 & $+$0.011 & .332 & .225 \\
\bottomrule
\end{tabular}
\end{table}

\subsection{Qwen2.5-7B}

\begin{table}[H]
\centering
\caption{Cross-dataset evaluation for Qwen2.5-7B: NormShift (with trick-question retry) vs.\ logit baselines.
         Trained on 50\% of ARC-Easy + TruthfulQA + ARC-Challenge,
         tested on held-out 50\% (in-domain) and unseen datasets (OOD).}
\label{tab:cross-dataset-qwen25-7b}
\small
\begin{tabular}{@{}llcccccc@{}}
\toprule
Dataset & Type & Acc & NS\,AUROC & MSP\,AUROC & $\Delta$ & NS\,HLCC & MSP\,HLCC \\
\midrule
Arc Easy             & in-dom & .853 & .716 & \textbf{.887} & $-$0.171 & 1.337 & 1.846 \\
Truthfulqa           & in-dom & .447 & \textbf{.706} & .620 & $+$0.086 & .358 & .304 \\
Arc Challenge        & in-dom & .793 & .636 & \textbf{.892} & $-$0.256 & 1.152 & 1.438 \\
Commonsenseqa        & OOD & .733 & .625 & \textbf{.797} & $-$0.172 & 1.019 & 1.297 \\
Openbookqa           & OOD & .667 & .641 & \textbf{.812} & $-$0.171 & .865 & 1.437 \\
Sciq                 & OOD & .847 & .523 & \textbf{.832} & $-$0.309 & 1.214 & 1.686 \\
Hellaswag            & OOD & .567 & .410 & \textbf{.742} & $-$0.332 & .402 & .735 \\
Mmlu                 & OOD & .493 & .568 & \textbf{.718} & $-$0.150 & .390 & .597 \\
\bottomrule
\end{tabular}
\end{table}

\subsection{Trick-Question Retry Summary}

\begin{table}[H]
\centering
\caption{NormShift-gated retry with trick-question prompt.
         When NormShift confidence $< 0.4$, the question is re-evaluated
         with a trick-question system prompt.  Rescue = wrong$\to$correct;
         Damage = correct$\to$wrong.}
\label{tab:retry-summary}
\small
\begin{tabular}{@{}llrrrr@{}}
\toprule
Model & Dataset & Retried & Rescues & Damage & Net \\
\midrule
Mistral-7B           & Arc Easy           & 13 & 0 & 0 & $+$0 \\
                     & Truthfulqa         & 78 & 25 & 0 & $+$25 \\
                     & Arc Challenge      & 28 & 3 & 1 & $+$2 \\
                     & Commonsenseqa      & 33 & 2 & 3 & $-$1 \\
                     & Openbookqa         & 18 & 1 & 0 & $+$1 \\
                     & Sciq               & 8 & 0 & 2 & $-$2 \\
                     & Hellaswag          & 40 & 0 & 2 & $-$2 \\
                     & Mmlu               & 95 & 3 & 3 & $+$0 \\
\midrule
Qwen2.5-7B           & Arc Easy           & 64 & 0 & 16 & $-$16 \\
                     & Truthfulqa         & 72 & 3 & 9 & $-$6 \\
                     & Arc Challenge      & 71 & 2 & 9 & $-$7 \\
                     & Commonsenseqa      & 76 & 2 & 13 & $-$11 \\
                     & Openbookqa         & 75 & 0 & 26 & $-$26 \\
                     & Sciq               & 89 & 1 & 12 & $-$11 \\
                     & Hellaswag          & 82 & 5 & 14 & $-$9 \\
                     & Mmlu               & 82 & 2 & 13 & $-$11 \\
\bottomrule
\end{tabular}
\end{table}



\section{Training Data Composition}
\label{app:training-data}

\begin{table}[H]
\centering
\caption{Training and evaluation data for NormShift confidence heads.
         All five models use identical data splits.}
\label{tab:data-composition}
\begin{tabular}{@{}llr@{}}
\toprule
Split & Dataset & Examples \\
\midrule
\multirow{3}{*}{Training (900 total)}
  & ARC-Easy       & 300 \\
  & TruthfulQA     & 300 \\
  & ARC-Challenge  & 300 \\
\midrule
Validation        & TruthfulQA & 100 \\
\midrule
Calibration eval  & TruthfulQA & 300 \\
\bottomrule
\end{tabular}
\end{table}

\paragraph{Dataset overlap.}
TruthfulQA appears in three roles: 300 training examples, 100 validation
examples (used for early stopping / checkpoint selection), and 300
calibration evaluation examples.  While we do not verify whether the
specific question instances overlap between these splits, the presence of
TruthfulQA in the training set means that calibration results on
TruthfulQA should be considered \emph{in-distribution} performance.

The strong AUROC (0.999) of NormShift on Mistral-7B TruthfulQA versus
its degradation on held-out datasets in \Cref{tab:cross-dataset}
directly reflects this overlap.  The 50/50 split evaluation
(\Cref{tab:cross-dataset-mistral-7b,tab:cross-dataset-qwen25-7b})
provides a cleaner test: NormShift still dominates in-domain for
Mistral-7B (AUROC 0.978 on held-out TruthfulQA) but not for Qwen2.5-7B
(AUROC 0.706), and degrades on OOD datasets for both models.

\paragraph{Checkpoint vs.\ calibration evaluation gap.}
A notable discrepancy exists between the checkpoint training metrics
and calibration evaluation.  For example, the Qwen2.5-7B combined head
achieves ECE\,=\,0.060 at checkpoint selection, but ECE\,=\,0.305 when
evaluated as a calibration baseline on TruthfulQA.  The Llama-3.1-8B
head shows a similar gap (checkpoint ECE\,=\,0.040 vs.\ calibration
eval ECE\,=\,0.235).  Only Mistral-7B shows close agreement (checkpoint
ECE\,=\,0.020 vs.\ calibration eval ECE\,=\,0.030).  This gap suggests
that the head overfits to its training distribution on some
architectures, and that checkpoint metrics alone are insufficient for
assessing generalisation.


\section{Methodology Notes}
\label{app:methodology-notes}

\paragraph{Post-hoc calibration data leakage.}
The post-hoc calibration methods (temperature scaling, isotonic
regression, Platt scaling) fit their calibration parameters on the same
300 TruthfulQA examples used for evaluation.  No held-out calibration
set is used.  As a result:
\begin{itemize}[nosep]
  \item Isotonic regression achieves ECE\,=\,0.000 by construction---it
        perfectly maps raw confidences to observed frequencies on the
        fitting data.  This does not reflect generalisation ability.
  \item Platt scaling achieves ECE\,=\,0.050, also partially an artefact
        of fitting on the evaluation set.
  \item Temperature scaling achieves ECE\,=\,0.515, which is \emph{worse}
        than uncalibrated MSP (0.353), because the learned temperature
        exacerbates overconfidence for this model.
\end{itemize}

These results should be interpreted as \emph{in-sample calibration
fits} rather than true generalisation performance.  A proper evaluation
would use a separate calibration holdout (e.g., 80/20 split of the 300
examples).

\paragraph{NormShift in-distribution advantage.}
The NormShift confidence head is trained using HLCC loss on data that
includes TruthfulQA examples.  On Mistral-7B, this produces
near-perfect discrimination (AUROC\,$\approx$\,1.0,
PRR\,$\approx$\,1.0), and the 50/50 split confirms strong in-domain
generalisation (AUROC 0.978 on held-out TruthfulQA, 0.930 on ARC-Easy,
0.919 on ARC-Challenge).  However, this strong performance does not
transfer to all architectures: on Qwen2.5-7B, NormShift achieves only
AUROC 0.706 on held-out TruthfulQA and loses to MSP on two of three
in-domain datasets.

On out-of-distribution data, AUROC drops to 0.49--0.80 (Mistral-7B)
and 0.41--0.64 (Qwen2.5-7B).  For Qwen2.5-7B, PRR metrics are
unreliable as a comparison since the retry mechanism changes the
accuracy.

This does not invalidate the norm-shift signal itself---the per-layer
standard deviation patterns do differ between correct and incorrect
predictions even on out-of-distribution data (\Cref{tab:layer-signals}).
Rather, it indicates that the \emph{learned decision boundary} is both
distribution-specific and architecture-dependent.  Mistral-7B's
normalisation dynamics appear to provide a particularly informative
basis for confidence estimation, while other architectures may require
different training strategies or larger training sets to achieve
comparable transfer.

\paragraph{Implications for general-purpose confidence.}
The current NormShift head should not be used as a general-purpose
confidence estimator for open-ended text generation without retraining on
data representative of the target domain.  For applications requiring
robust confidence estimates across diverse inputs, logit-based methods
(particularly MSP or entropy) remain more reliable due to their implicit
adaptation to the model's output distribution.


% ============================================================
% BIBLIOGRAPHY
% ============================================================
\bibliographystyle{plainnat}
\bibliography{hlcc_paper}

\end{document}
